\documentclass[12pt,a4paper,titlepage]{article}

% Mathematics
\usepackage{amsmath,amsthm,amssymb,amsfonts}
\usepackage{mathtools}

% Layout (asymmetric margins: 4cm left for binding, 2cm right)
\usepackage[a4paper,lmargin=4cm,rmargin=2cm,tmargin=2.5cm,bmargin=2.5cm]{geometry}
\linespread{1.25}

% Headers and footers
\usepackage{fancyhdr}
\pagestyle{fancy}
\fancyhf{}
\fancyhead[L]{\nouppercase{\leftmark}}
\fancyhead[R]{\thepage}
\fancyfoot[C]{\footnotesize Numerical Investigation of RH via DQPTs}
\renewcommand{\headrulewidth}{0.4pt}
\renewcommand{\footrulewidth}{0.4pt}
\fancypagestyle{plain}{%
  \fancyhf{}
  \fancyhead[R]{\thepage}
  \fancyfoot[C]{\footnotesize Numerical Investigation of RH via DQPTs}
  \renewcommand{\headrulewidth}{0pt}
  \renewcommand{\footrulewidth}{0.4pt}
}

% References
\usepackage[numbers,sort&compress]{natbib}
\usepackage{url}
\usepackage{hyperref}
\hypersetup{colorlinks=true,linkcolor=blue,citecolor=blue,urlcolor=blue}

% Tables and figures
\usepackage{graphicx}
\usepackage{float}
\usepackage{booktabs}
\usepackage{longtable}
\usepackage{caption}
\captionsetup{font=small,labelfont=bf}

% Code listings
\usepackage{listings}
\lstset{basicstyle=\ttfamily\small,breaklines=true,frame=single}

% Theorem environments
\newtheorem{theorem}{Theorem}[section]
\newtheorem{lemma}[theorem]{Lemma}
\newtheorem{proposition}[theorem]{Proposition}
\newtheorem{corollary}[theorem]{Corollary}
\newtheorem{definition}[theorem]{Definition}
\theoremstyle{remark}
\newtheorem{remark}[theorem]{Remark}

% Section-based equation numbering: (1.1), (1.2), (2.1), ...
\numberwithin{equation}{section}

% Notation
\DeclareMathOperator{\re}{Re}
\DeclareMathOperator{\im}{Im}
\DeclareMathOperator{\Tr}{Tr}
\newcommand{\Z}{\mathbb{Z}}
\newcommand{\R}{\mathbb{R}}
\newcommand{\C}{\mathbb{C}}
\newcommand{\N}{\mathbb{N}}
\newcommand{\eps}{\varepsilon}

\begin{document}

%% ---- Title page ----
\begin{titlepage}
  \centering
  \vspace*{3cm}
  {\LARGE\bfseries
    Numerical Investigation of the Riemann Hypothesis\\[0.4cm]
    via Dynamical Quantum Phase Transitions\\[0.4cm]
  }
  {\Large Implementation Plan and Theoretical Framework}

  \vspace{2cm}
  {\large \today}

  \vspace{1.5cm}
  {\large\url{https://github.com/Protosophos/structure-space}}

  \vfill
\end{titlepage}

%% ---- Abstract (page ii, Roman numerals) ----
\pagenumbering{Roman}

\begin{abstract}
We present a detailed plan for the first public classical-computer implementation of
the framework described by Wei et~al.~\cite{wei2025}, which establishes an exact
equivalence between the Riemann Hypothesis~(RH) and the confinement of dynamical
quantum phase transitions~(DQPTs) to the critical line $\re(s) = \tfrac{1}{2}$ in
two carefully constructed quantum systems. Our implementation classically simulates
both quantum systems---a probe spin coupled to a bath with logarithmic energy spectrum
and a generalized Loschmidt echo construction---computing their respective order
parameters on a fine grid in the $(\beta, t)$ plane. The software is written in C++
with CUDA acceleration for large-scale parameter sweeps. By comparing detected
phase-transition points against the known nontrivial zeros of the Riemann zeta function
and scanning for off-critical-line zeros, we provide independent numerical evidence
bearing on the Riemann Hypothesis through the lens of quantum many-body physics.

We emphasize that this numerical approach cannot prove the Riemann Hypothesis; it can,
however, either discover a counterexample (disproving~RH) or provide additional
computational evidence consistent with~RH, while simultaneously validating the
Wei et~al.\ theoretical framework~\cite{wei2025} against established zero tables.

\medskip\noindent
\textbf{Notation.} Throughout this document, $\log$ denotes the natural logarithm
(base~$e$). We write $\log_2$ explicitly when base-2 logarithm is intended.
\end{abstract}

%% ---- Table of contents (Roman numerals) ----
\newpage
\tableofcontents

%% ---- Main content (Arabic numerals, starting at 1) ----
\newpage
\pagenumbering{arabic}
\setcounter{page}{1}

%% ====================================================================
\section{Introduction and Motivation}
\label{sec:intro}
%% ====================================================================

\subsection{The Riemann Hypothesis}
\label{sec:rh}

The Riemann zeta function is defined for $\re(s) > 1$ by the absolutely convergent
series
\begin{equation}\label{eq:zeta-def}
  \zeta(s) = \sum_{n=1}^{\infty} n^{-s},
\end{equation}
and extended to the entire complex plane (except for a simple pole at $s=1$) by
analytic continuation. The functional equation
\begin{equation}\label{eq:func-eq}
  \zeta(s) = \chi(s)\,\zeta(1-s),
  \qquad
  \chi(s) = 2^s \pi^{s-1} \sin\!\bigl(\tfrac{\pi s}{2}\bigr)\,\Gamma(1-s),
\end{equation}
equivalently expressed using the completed zeta function
\begin{equation}\label{eq:xi-def}
  \xi(s) = \tfrac{1}{2}\,s(s-1)\,\pi^{-s/2}\,\Gamma\!\bigl(\tfrac{s}{2}\bigr)\,\zeta(s),
\end{equation}
which satisfies the symmetric form $\xi(s) = \xi(1-s)$, implies that $\zeta(s)$ vanishes
at $s = -2, -4, -6, \ldots$ (the \emph{trivial zeros}). The Riemann
Hypothesis~\cite{riemann1859} asserts:

\medskip
\noindent\textbf{RH.} \emph{All nontrivial zeros of $\zeta(s)$ have real part equal
to $\tfrac{1}{2}$.}
\medskip

\noindent
Equivalently, if $\zeta(\sigma + it) = 0$ with $0 < \sigma < 1$ and $t \in \R$, then
$\sigma = \tfrac{1}{2}$. This remains the most important open problem in mathematics,
with profound implications for the distribution of prime numbers, the error term in the
prime counting function, and numerous results across number theory, analysis, and
mathematical physics that are conditional on~RH.

\subsection{Computational Approaches to RH}
\label{sec:computational}

The dominant computational approach to~RH has been direct numerical verification:
computing $\zeta(\tfrac{1}{2} + it)$ along the critical line and certifying that all
zeros up to some height~$T$ lie on it. Odlyzko~\cite{odlyzko1989,odlyzko2001} computed
specific high zeros of~$\zeta$---around the $10^{20}$th and $10^{22}$nd zeros---and
large blocks of consecutive zeros near these heights, providing detailed statistical
data on zero spacing and pair correlation. Gourdon~\cite{gourdon2004} systematically
verified that all of the first $10^{13}$ nontrivial zeros lie on the critical line,
using the Odlyzko--Sch\"onhage algorithm based on the Riemann--Siegel formula. These
computations exploit highly optimized algorithms specific to the zeta function and
achieve extraordinary precision.

Our approach differs fundamentally. Rather than computing $\zeta$ directly, we
\emph{classically} evaluate the mathematical expressions (partial Dirichlet sums,
Loschmidt amplitudes, rate functions) that arise from quantum-mechanical systems
whose observable phase transitions correspond exactly to the zeros
of~$\zeta$.\footnote{To be precise: no quantum hardware or quantum simulation is
involved. The computation reduces to evaluating partial sums of Dirichlet series and
related functions on a classical computer. The quantum-mechanical language (Hamiltonians,
coherence functions, Loschmidt echoes) provides the \emph{framework} and
\emph{interpretation}, not the computational method. The value of this framing lies in
connecting the Riemann Hypothesis to the well-developed theory of dynamical quantum phase
transitions, not in any computational advantage over direct zeta evaluation.} This
provides:
\begin{enumerate}
  \item An independent computational pathway to zero detection, serving as a cross-check.
  \item Physical intuition about the structure of zeta zeros through the language of
        quantum phase transitions.
  \item A connection (though not a direct realization) to the Hilbert--P\'olya
        program\footnote{The Hilbert--P\'olya conjecture, traditionally attributed to
        c.~1914 (though the earliest written record is later), posits the existence of
        a self-adjoint operator whose eigenvalues are the imaginary parts of the
        nontrivial zeta zeros. If such an operator exists, the reality of its
        eigenvalues would immediately imply~RH.}, linking~RH to quantum dynamical
        phenomena rather than to a single spectral operator.
\end{enumerate}

\subsection{The Hilbert--P\'olya Conjecture and DQPTs}
\label{sec:hilbert-polya}

Berry and Keating~\cite{berrykeating1999} proposed that the Hilbert--P\'olya operator
might be related to the quantization of the classical Hamiltonian $H = xp$, connecting~RH
to quantum chaos.

The Wei et~al.\ construction~\cite{wei2025} takes a different but related path. Instead
of seeking a single operator whose spectrum encodes the zeros, they construct quantum
systems whose \emph{dynamical} phase transitions---non-analyticities in a rate function
as a function of real time---occur precisely at the zeta zeros. The key insight is that
the Loschmidt amplitude of these systems, when evaluated at inverse temperature~$\beta$,
is proportional to the Dirichlet eta function $\eta(\beta + it)$ (System~1) or the
Hardy $Z$-function (System~2). The~RH is then equivalent to the statement that DQPTs
occur only at $\beta = \tfrac{1}{2}$.

This reformulation shifts~RH from a statement about complex analysis to one about
quantum statistical mechanics, opening new avenues for both theoretical and numerical
investigation.

\subsection{Scope and Limitations of the DQPT Reformulation}
\label{sec:scope}

We emphasize an important conceptual point: the DQPT--RH equivalence established by
Wei et~al.\ is a \emph{reformulation} of the Riemann Hypothesis in the language of
quantum statistical mechanics, not a proof mechanism. The equivalence states that~RH
holds if and only if DQPTs are confined to $\beta = \tfrac{1}{2}$, but this
reformulation does not by itself provide any new leverage toward proving (or
disproving)~RH. The difficulty of the problem is preserved: it is merely re-expressed
in physical language.

In particular:
\begin{enumerate}
  \item The quantum systems (System~1 and System~2) are \emph{constructed} so that
        their partition functions and order parameters reproduce the Dirichlet eta
        function and Hardy $Z$-function. The physics does not generate~RH as a
        consequence of some independent physical principle; rather, the mathematics
        of~RH is \emph{encoded} into the physics by design.
  \item The numerical investigation described in this document can provide computational
        evidence consistent with (or inconsistent with)~RH, but it cannot prove~RH.
        This is a fundamental limitation shared by all numerical approaches.
  \item The value of the DQPT framework lies in (a)~providing physical intuition about
        the structure of zeta zeros, (b)~connecting~RH to the well-developed theory of
        quantum phase transitions, and (c)~potentially suggesting new analytical
        approaches inspired by the physics (see Section~\ref{sec:directions}).
\end{enumerate}

%% ====================================================================
\section{Theoretical Framework}
\label{sec:theory}
%% ====================================================================

\subsection{Dynamical Quantum Phase Transitions}
\label{sec:dqpt}

Dynamical quantum phase transitions (DQPTs) were introduced by Heyl, Polkovnikov, and
Kehrein~\cite{heyl2013} as real-time analogs of equilibrium phase transitions. The
central object is the Loschmidt amplitude:

\begin{definition}[Loschmidt amplitude]\label{def:loschmidt}
Given an initial state $|\psi_0\rangle$ (typically the ground state of a pre-quench
Hamiltonian~$H_0$) and a post-quench Hamiltonian~$H$, the Loschmidt amplitude is
\begin{equation}\label{eq:loschmidt}
  \mathcal{G}(t) = \langle\psi_0| e^{-iHt} |\psi_0\rangle.
\end{equation}
The return rate (or rate function) is defined as
\begin{equation}\label{eq:rate-general}
  g(t) = -\frac{1}{N}\ln|\mathcal{G}(t)|^2,
\end{equation}
where $N$ is the system size. A DQPT occurs at time~$t^*$ if $g(t)$ is non-analytic
at~$t = t^*$.
\end{definition}

The connection to equilibrium statistical mechanics arises through the formal
correspondence $t \to -i\beta$ (Wick rotation), under which the Loschmidt amplitude
becomes
\begin{equation}\label{eq:wick}
  \mathcal{G}(-i\beta) = \langle\psi_0| e^{-\beta H} |\psi_0\rangle.
\end{equation}
For the specific case where $|\psi_0\rangle$ is a uniform superposition over energy
eigenstates, this equals $Z(\beta)/\dim(\mathcal{H})$. In general, the relation depends
on the overlap of $|\psi_0\rangle$ with the eigenstates of~$H$.

Non-analyticities in the free energy density $f(\beta) = -(1/N)\ln Z(\beta)$ as a
function of inverse temperature correspond to equilibrium phase transitions. By the
Lee--Yang theorem~\cite{leeyang1952}, these non-analyticities are controlled by the
zeros of the partition function in the complex $\beta$-plane. Analogously, DQPTs are
controlled by the zeros of the Loschmidt amplitude in the complex time plane. When
these zeros cross the real time axis, the rate function $g(t)$ exhibits non-analytic
behavior (cusps, divergences, or discontinuities in derivatives), signaling a
DQPT~\cite{heyl2013}.

\paragraph{Connection to Lee--Yang zeros.}
In the Lee--Yang theory~\cite{leeyang1952,fisher1965}, equilibrium phase transitions
arise when zeros of the partition function $Z(\beta,h)$---viewed as a function of
complex fugacity or magnetic field---pinch the real axis in the thermodynamic limit.
The density of these zeros determines the order and nature of the transition. In the
DQPT framework, the analogous role is played by the Fisher zeros of the Loschmidt
amplitude in the complex time plane. The Wei et~al.\ construction arranges for these
Fisher zeros to coincide with the nontrivial zeros of the Riemann zeta function.


\subsection{The Wei et al.\ Construction}
\label{sec:wei}

\subsubsection{System~1: Probe Spin Coherence with Logarithmic Spectrum}
\label{sec:system1}

\paragraph{Setup.}
Consider a composite system consisting of a probe spin-$\tfrac{1}{2}$ coupled to a
bath of $N$ quantum levels. The bath Hamiltonian has a logarithmic energy spectrum:
\begin{equation}\label{eq:H0}
  H_0 = \sum_{n=1}^{N} \ln(n)\,|n\rangle\langle n|,
\end{equation}
where $\{|n\rangle\}_{n=1}^{N}$ is an orthonormal basis for the bath Hilbert space.

\begin{remark}[On the spectrum]\label{rem:spectrum}
The logarithmic spectrum $E_n = \ln(n)$ is reverse-engineered to produce the Dirichlet
series in the partition function; it is not a spectrum that arises naturally in a
standard condensed-matter or atomic physics context. However, logarithmic spectra do
arise in the theory of quantum mechanics on hyperbolic surfaces: the Selberg zeta
function
\begin{equation}\label{eq:selberg}
  Z_\Gamma(s) = \prod_{\{p\}} \prod_{k=0}^{\infty}
  \bigl(1 - e^{-(s+k)\ell_p}\bigr),
\end{equation}
associated to a discrete subgroup $\Gamma$ of $\mathrm{SL}(2,\R)$, encodes lengths
$\ell_p$ of primitive closed geodesics, and the corresponding quantum Hamiltonian (the
Laplace--Beltrami operator on the quotient surface) has a spectral theory intimately
related to zeta functions~\cite{edwards1974}. This provides a potential natural
realization of the logarithmic spectrum, though the precise connection to the
Wei et~al.\ construction remains to be established.
\end{remark}

The probe spin couples to the bath such that the effective Hamiltonian for the coupled
evolution is
\begin{equation}\label{eq:H-coupled}
  H = \sigma_z \otimes H_0,
\end{equation}
where $\sigma_z$ is the Pauli-$Z$ operator acting on the probe spin.

\paragraph{Thermal state preparation.}
The bath is prepared in a thermal (Gibbs) state at inverse temperature~$\beta$:
\begin{equation}\label{eq:thermal}
  \rho_{\mathrm{bath}}(\beta) = \frac{e^{-\beta H_0}}{Z(\beta,H_0)},
\end{equation}
where the partition function is
\begin{equation}\label{eq:partition}
  Z(\beta,H_0) = \Tr\bigl[e^{-\beta H_0}\bigr]
  = \sum_{n=1}^{N} e^{-\beta\ln(n)}
  = \sum_{n=1}^{N} n^{-\beta}.
\end{equation}
Note that $Z(\beta,H_0)$ is precisely the partial zeta function. In the thermodynamic
limit $N \to \infty$ and for $\beta > 1$, this converges to
$\zeta(\beta)$~\cite{titchmarsh1986}.

The probe spin is prepared in the state
$|{+}\rangle = (|{\uparrow}\rangle + |{\downarrow}\rangle)/\sqrt{2}$.

\paragraph{Time evolution.}
The probe-bath system is prepared in the product state
\begin{equation}\label{eq:initial-state}
  \rho_{\mathrm{total}} = |{+}\rangle\langle{+}| \otimes \rho_{\mathrm{bath}}(\beta),
\end{equation}
where $|{+}\rangle\langle{+}|$ is the pure state projector for the probe spin and
$\rho_{\mathrm{bath}}(\beta)$ is the thermal state defined above. Under the coupled
Hamiltonian $H = \sigma_z \otimes H_0$, the total state evolves as
\begin{equation}\label{eq:evolution}
  \rho_{\mathrm{total}}(t) = e^{-iHt}\,\rho_{\mathrm{total}}\,e^{iHt}.
\end{equation}
The reduced density matrix of the probe spin, obtained by tracing over the bath
(partial trace), yields the coherence (off-diagonal element):
\begin{equation}\label{eq:coherence}
  \langle\sigma_+\rangle(t) = \Tr_{\mathrm{bath}}\bigl[\rho_{\mathrm{bath}}\,e^{2iH_0 t}\bigr].
\end{equation}

\paragraph{Derivation of the accumulated phase factor.}
We compute:
\begin{align}\label{eq:coherence-deriv}
  \langle\sigma_+\rangle(t)
  &= \Tr\bigl[\rho_{\mathrm{bath}}\,e^{2iH_0 t}\bigr] \nonumber\\
  &= \frac{1}{Z(\beta)} \sum_{n=1}^{N} e^{-\beta\ln(n)}\,e^{2i\ln(n)\,t} \nonumber\\
  &= \frac{1}{Z(\beta)} \sum_{n=1}^{N} n^{-\beta}\,n^{2it} \nonumber\\
  &= \frac{1}{Z(\beta)} \sum_{n=1}^{N} n^{-(\beta - 2it)} \nonumber\\
  &= \frac{S_N(\beta - 2it)}{Z(\beta)},
\end{align}
where $S_N(s) = \sum_{n=1}^{N} n^{-s}$ is the partial Dirichlet series.

Note that the coherence $\langle\sigma_+\rangle$ involves $S_N(\beta - 2it)$, not
$S_N(\beta + it)$. Wei et~al.~\cite{wei2025} introduce three modifications to obtain
the accumulated phase factor $\mathcal{L}(\beta,t)$:
\begin{enumerate}
  \item[(a)] A reparametrization $t \to t/2$ absorbs the factor of~2: replacing $2t$
    by~$t$ in the exponent gives $S_N(\beta - it)$. This is purely a rescaling of the
    time variable and does not affect the location of zeros (a zero of $S_N(\beta - it_0)$
    at time~$t_0$ corresponds to a zero of $S_N(\beta - 2it_0')$ at time $t_0' = t_0/2$).
  \item[(b)] A sign convention: the coherence involves $n^{-(\beta - it)} =
    n^{-\beta}\,n^{it}$, while the standard Dirichlet series uses
    $n^{-(\beta + it)} = n^{-\beta}\,n^{-it}$. The passage from $-it$ to $+it$
    corresponds to complex conjugation of the time evolution phase, equivalent to
    choosing the opposite sign convention for the coupling Hamiltonian. This does not
    affect the zero locations (if $\eta(\beta + it_0) = 0$, then $\eta(\beta - it_0) = 0$
    by the reflection principle, since the eta coefficients are real).
  \item[(c)] An alternating sign $(-1)^{n+1}$ is introduced, corresponding physically to
    a staggered coupling between the probe and even/odd bath levels. This converts the
    ordinary Dirichlet series into the alternating Dirichlet eta series.\footnote{This
    staggered coupling requires that the probe spin distinguishes the parity of the
    energy index~$n$ and reverses its coupling accordingly. Since $n$ is an energy index
    rather than a spatial coordinate, this coupling is \emph{non-local in energy space}
    and does not arise naturally in any known many-body system. Together with the
    logarithmic spectrum $E_n = \ln n$, this reinforces the point made in
    Section~\ref{sec:scope}: the quantum system is \emph{engineered} to reproduce the
    Dirichlet eta function, and the ``physics'' is a mathematical reformulation rather
    than an independent physical mechanism.}
\end{enumerate}
With these three modifications, the accumulated phase factor becomes:
\begin{equation}\label{eq:L-def}
  \mathcal{L}(\beta,t) = \frac{1}{Z(\beta)}
  \sum_{n=1}^{N} (-1)^{n+1}\,n^{-(\beta+it)}.
\end{equation}
The factor $(-1)^{n+1}$ converts the ordinary Dirichlet series into the alternating
Dirichlet series, which defines the Dirichlet eta function:
\begin{equation}\label{eq:eta-def}
  \eta(s) = \sum_{n=1}^{\infty} (-1)^{n+1}\,n^{-s}
  = (1 - 2^{1-s})\,\zeta(s).
\end{equation}

\paragraph{Why $\mathcal{L}$ vanishes at zeta zeros.}
We now prove the key result. In the thermodynamic limit $N \to \infty$:

\medskip\noindent
\textbf{Step~1.} The numerator of $\mathcal{L}(\beta,t)$ converges:
\begin{equation}\label{eq:eta-convergence}
  \sum_{n=1}^{\infty} (-1)^{n+1}\,n^{-(\beta+it)} = \eta(\beta+it).
\end{equation}
This series converges for $\re(\beta+it) = \beta > 0$ by Dirichlet's test (Abel
summation): the partial sums of $(-1)^{n+1}$ are bounded and $n^{-\beta}$ decreases
monotonically to zero~\cite{titchmarsh1986}. (Note: the classical Leibniz alternating
series test applies only to real series; for complex $s = \beta + it$ with $t \ne 0$,
convergence follows from Dirichlet's test.)

\medskip\noindent
\textbf{Step~2.} The Dirichlet eta function satisfies the identity
\begin{equation}\label{eq:eta-zeta}
  \eta(s) = (1 - 2^{1-s})\,\zeta(s)
\end{equation}
for all $s \in \C$ (by analytic continuation). This identity is proved by observing
that for $\re(s) > 1$:
\begin{equation}\label{eq:eta-zeta-proof}
  \zeta(s) - \eta(s) = 2\sum_{n=1}^{\infty}(2n)^{-s} = 2^{1-s}\,\zeta(s),
\end{equation}
whence $\eta(s) = \zeta(s) - 2^{1-s}\zeta(s) = (1-2^{1-s})\zeta(s)$. Both sides are
meromorphic on~$\C$ and agree on $\re(s) > 1$; by the identity theorem for meromorphic
functions, they agree everywhere on~$\C$.

\medskip\noindent
\textbf{Step~3.} The prefactor $(1 - 2^{1-s})$ vanishes when $2^{1-s} = 1$, i.e., when
$(1-s)\ln 2 = 2\pi ik$ for some integer~$k$. Writing $s = \beta + it$:
\begin{equation}\label{eq:prefactor-eq}
  (1 - \beta - it)\ln 2 = 2\pi ik.
\end{equation}
Separating real and imaginary parts:
\begin{align}
  \text{Real:}\quad & (1-\beta)\ln 2 = 0, \quad\text{hence } \beta = 1. \label{eq:pf-real}\\
  \text{Imaginary:}\quad & -t\ln 2 = 2\pi k, \quad\text{hence }
    t = -2\pi k/\ln 2 = 2\pi(-k)/\ln 2. \label{eq:pf-imag}
\end{align}
Since $k$ ranges over all integers, the set of solutions is $t = 2\pi m/\ln 2$ for
$m \in \Z$ (relabeling $m = -k$). The prefactor zeros thus lie at
$s = 1 + i(2\pi m/\ln 2)$ for $m \in \Z$. Since $\beta = 1$ lies \emph{outside} the
open critical strip $0 < \beta < 1$ (and no nontrivial zeta zero has $\re(s) = 1$
by the prime number theorem~\cite{titchmarsh1986}), these prefactor zeros never
coincide with nontrivial zeta zeros.

\medskip\noindent
\textbf{Step~4.} Therefore, for $s = \beta + it$ with $0 < \beta < 1$:
\begin{equation}\label{eq:eta-zero-equiv}
  \eta(s) = 0 \;\Longleftrightarrow\; \zeta(s) = 0.
\end{equation}

\medskip\noindent
\textbf{Step~5: The partition function divergence.}
The partition function $Z(\beta) = \sum_{n=1}^{N} n^{-\beta}$ converges to
$\zeta(\beta)$ only for $\beta > 1$. For $\beta = \tfrac{1}{2}$ (the critical case
for~RH), $Z(\tfrac{1}{2}) = \sum 1/\sqrt{n}$ diverges as $2\sqrt{N}$ by the integral
test~\cite{titchmarsh1986}.

This means the naive ratio $\mathcal{L} = \eta / Z$ is a $0/\infty$ form at the zeta
zeros when $\beta = \tfrac{1}{2}$. We must handle this carefully.

Following Wei et~al.~\cite{wei2025}, the physically meaningful quantity is the
\emph{finite-$N$} expression:
\begin{equation}\label{eq:LN-def}
  \mathcal{L}_N(\beta,t) = \frac{\sum_{n=1}^{N}(-1)^{n+1}n^{-(\beta+it)}}
  {\sum_{n=1}^{N} n^{-\beta}}.
\end{equation}
For finite~$N$, both numerator and denominator are finite. The key observation is that
the \emph{rate function} (free energy density) has a well-defined thermodynamic limit:
\begin{equation}\label{eq:F1-def}
  \mathcal{F}_1(\beta,t) = -\lim_{N\to\infty} \frac{1}{\ln N}\,
  \ln|\mathcal{L}_N(\beta,t)|.
\end{equation}
To see this, note that for the numerator:
\begin{itemize}
  \item At a zeta zero $s_0 = \beta_0 + it_0$: the partial eta sum
    $S_N^{\eta}(s_0) = O(N^{-\beta_0})$ by the Dirichlet test remainder
    bound~\cite{titchmarsh1986}.
  \item Away from zeros: $S_N^{\eta}(s) \to \eta(s) \ne 0$, a nonzero constant.
\end{itemize}
For the denominator $Z_N(\beta)$:
\begin{itemize}
  \item For $0 < \beta < 1$: $Z_N(\beta) \sim N^{1-\beta}/(1-\beta)$ by
    Euler--Maclaurin summation~\cite{titchmarsh1986}.
\end{itemize}
Therefore:
\begin{equation}\label{eq:LN-scaling}
  |\mathcal{L}_N(\beta,t)| \sim |S_N^{\eta}(s)|\,\frac{1-\beta}{N^{1-\beta}}.
\end{equation}
At a zero ($s = s_0$): $|S_N^{\eta}(s_0)| = O(N^{-\beta_0})$, so
$|\mathcal{L}_N| \sim N^{\beta-\beta_0-1}$.
Away from a zero: $|S_N^{\eta}(s)| = O(1)$, so $|\mathcal{L}_N| \sim N^{\beta-1}$.

We verify the limit~\eqref{eq:F1-def} exists. Away from zeros,
$\ln|\mathcal{L}_N| = \ln|\eta(s)(1-\beta)| - (1-\beta)\ln N$, giving:
\begin{align}
  \mathcal{F}_1 &= -\lim_{N\to\infty}\frac{1}{\ln N}
    \bigl[\ln|\eta(s)(1-\beta)| - (1-\beta)\ln N\bigr] \nonumber\\
  &= -\lim_{N\to\infty}\frac{\ln|\eta(s)(1-\beta)|}{\ln N} + (1-\beta)
    = 1-\beta. \label{eq:F1-away}
\end{align}
At a zero, $\ln|\mathcal{L}_N| = \ln D + (\beta-\beta_0-1)\ln N$:
\begin{align}
  \mathcal{F}_1 &= -\lim_{N\to\infty}\frac{1}{\ln N}
    \bigl[\ln D + (\beta-\beta_0-1)\ln N\bigr] \nonumber\\
  &= \beta_0 + 1 - \beta. \label{eq:F1-zero}
\end{align}
Both limits exist and are finite. This confirms that $\mathcal{F}_1$ is well-defined
despite the divergence of $Z(\beta)$ for $\beta \le 1$.

\medskip\noindent
\textbf{Step~6: Equivalence statement.} The Riemann Hypothesis is equivalent to:
\begin{equation}\label{eq:rh-dqpt-equiv}
  \text{For all } t \in \R:\;
  \mathcal{F}_1(\beta,t) \text{ exhibits a singularity (DQPT) only at }
  \beta = \tfrac{1}{2}.
\end{equation}
A counterexample to~RH---a zero $\rho = \beta_0 + it_0$ with
$\beta_0 \ne \tfrac{1}{2}$---would manifest as a singularity in $\mathcal{F}_1$ at
$\beta = \beta_0$. This is numerically detectable by scanning $\mathcal{F}_1(\beta,t)$
over a grid in the $(\beta,t)$ plane and looking for anomalous values of
$\mathcal{F}_1$ at $\beta \ne \tfrac{1}{2}$.

\paragraph{Convergence of the eta partial sums.}
The Dirichlet eta series converges for $\re(s) > 0$ by Dirichlet's test (Abel
summation), but only conditionally (not absolutely) for $0 < \re(s) \le 1$. The
convergence rate is:
\begin{equation}\label{eq:eta-remainder}
  |\eta(s) - \eta_N(s)| \le C(s)\,(N+1)^{-\beta},
\end{equation}
where $C(s) = O(1)$ for $s$ bounded away from $\re(s) = 0$, but $C(s)$ grows as
$|\im(s)|$ increases for fixed $\re(s)$ near~0. This is $O(N^{-1/2})$ at
$\beta = \tfrac{1}{2}$, requiring $N \sim 10^{2k}$ to achieve $k$ digits of precision.
For double precision (${\sim}15$ digits), we need $N \sim 10^{30}$, which is infeasible
by direct summation. We address acceleration techniques in Section~\ref{sec:methods}.

\paragraph{Numerical acceleration.}
To circumvent the slow convergence, we use:
\begin{enumerate}
  \item[(a)] The Borwein--Cohen--Villegas--Zagier acceleration~\cite{borwein2000}:
    transforms the conditionally convergent eta series into a geometrically convergent
    series with error $O(3^{-N})$, requiring only $N \sim 40$ terms for double precision.
  \item[(b)] The Euler--Maclaurin summation formula applied to the non-alternating sum
    $S_N(s)$, combined with the identity $\eta(s) = (1-2^{1-s})\zeta(s)$ and the
    Riemann--Siegel formula for~$\zeta$.
\end{enumerate}
Both methods are detailed in Section~\ref{sec:methods}.

\begin{remark}[Trivial zeros]\label{rem:trivial}
The trivial zeros of $\zeta(s)$ at $s = -2,-4,-6,\ldots$ lie outside the critical strip
(they have $\re(s) < 0$). Since the eta series only converges for $\re(s) > 0$, and we
restrict our analysis to $0 < \beta < 1$, the trivial zeros do not appear in our
framework. Furthermore, the prefactor $(1-2^{1-s})$ does not vanish at the trivial zeros
(since $1-s = 3,5,7,\ldots$ gives $2^{1-s} = 2^3,2^5,\ldots \ne 1$).
\end{remark}

\paragraph{Why zeros of $\mathcal{L}$ correspond to DQPTs.}
A DQPT is defined as a non-analyticity (cusp, discontinuity, or divergence) in the rate
function $\mathcal{F}(\beta,t)$, not merely a zero of the order parameter
$\mathcal{L}(\beta,t)$. We show these are equivalent in this framework.

Consider the rate function $\mathcal{F}_1(\beta,t) = -(1/\ln N)\ln|\mathcal{L}_N(\beta,t)|$
as $N \to \infty$:
\begin{enumerate}
  \item[(i)] For $(\beta,t)$ such that $\eta(\beta+it) \ne 0$:
    the numerator $\eta_N(s) \to \eta(s) \ne 0$, a bounded nonzero analytic function
    of~$s$. Write $\eta_N(s) = \eta(s) + O(N^{-\beta})$ by the Dirichlet test remainder
    bound. The denominator $Z_N(\beta) = N^{1-\beta}/(1-\beta) + O(1)$ by
    Euler--Maclaurin~\cite{titchmarsh1986}. Hence
    \begin{equation}
      |\mathcal{L}_N| = |\eta(s)|\,(1-\beta)\,N^{\beta-1}\,
      \bigl(1 + O(N^{-\min(\beta,1-\beta)})\bigr).
    \end{equation}
    The leading behavior gives $\mathcal{F}_1 \to (1-\beta)$, a smooth function.
  \item[(ii)] At a zero $s_0 = \beta_0 + it_0$ of $\eta$ (equivalently, of $\zeta$):
    write $\eta(s) = \eta'(s_0)(s-s_0) + O(|s-s_0|^2)$ near~$s_0$ (assuming simple
    zero\footnote{For zeros of multiplicity~$m$, see Remark~\ref{rem:mult}.}). At
    $s = s_0$ exactly: $\eta_N(s_0) = O(N^{-\beta_0})$ by the Dirichlet test remainder
    bound (since the full series converges to $\eta(s_0) = 0$). Hence
    $|\mathcal{L}_N(\beta_0,t_0)| = O(N^{\beta-\beta_0-1})$ and
    $\mathcal{F}_1 \to (\beta_0+1-\beta)$.
\end{enumerate}
The rate function jumps discontinuously from $(1-\beta)$ to $(\beta_0+1-\beta)$ at the
zero, a discontinuity of magnitude~$\beta_0$. Since $\beta_0 > 0$ for all nontrivial
zeros, this is a genuine non-analyticity---a first-order DQPT in the classification of
Heyl et~al.~\cite{heyl2013}.

More precisely, the derivative $d\mathcal{F}_1/dt$ diverges at $t = t_0$ (the rate
function has a cusp), because $|\mathcal{L}_N|$ passes through a minimum of order
$N^{\beta-\beta_0-1}$ and the logarithm amplifies this into a peak of height
$O(\ln N)$ in $\mathcal{F}_1$. In the thermodynamic limit, this becomes a genuine
singularity.

Therefore: zeros of $\mathcal{L}$ correspond one-to-one with DQPTs (singularities of
$\mathcal{F}_1$). This completes the heuristic argument.

\begin{remark}[Rigor]\label{rem:rigor}
The above argument is carried out at the level of formal asymptotics. A fully rigorous
statement requires establishing (i)~the existence of the thermodynamic limit
$N \to \infty$ for $\mathcal{F}_1$, with uniform convergence in $(\beta,t)$ on compact
subsets of $(0,1) \times \R$ away from zeros, and (ii)~that the singularity structure
persists in the limit. We verify both conditions numerically (Section~\ref{sec:validation}).
\end{remark}

\begin{remark}[Lehmer phenomenon]\label{rem:lehmer}
Occasionally, two zeta zeros lie extremely close together. A notable early example is
the pair near $t \approx 7005$ (gap approximately~0.04), identified by
Lehmer~\cite{lehmer1956}. Much smaller gaps ($< 10^{-3}$) are known at larger
heights~\cite{odlyzko1989}. Our zero-detection algorithm (Section~\ref{sec:zero-detection})
accounts for this with adaptive $t$-refinement.
\end{remark}

\paragraph{Catastrophic cancellation at large $t$.}
For large~$t$, the summand $n^{-(\beta+it)} = n^{-\beta}e^{-it\ln n}$ involves rapidly
oscillating phases. The partial sum $S_N(s) = \sum_{n=1}^{N} n^{-s}$ exhibits extensive
cancellation. By the variance heuristic (treating phases as approximately random), the
expected magnitude is $O\bigl((\sum_{n=1}^{N} n^{-2\beta})^{1/2}\bigr)$. For
$\beta = \tfrac{1}{2}$, this gives $O(\sqrt{\ln N})$; for $\beta > \tfrac{1}{2}$,
$O(1)$; for $\beta < \tfrac{1}{2}$, $O(N^{1/2-\beta})$. More rigorous bounds follow
from van der Corput estimates~\cite{titchmarsh1986}:
$|S_N(\beta+it)| = O(N^{1-\beta}t^{-1/2} + t^{1/2})$ for $t > 0$.

The practical consequence: with $N$ terms each contributing $O(\eps_{\mathrm{mach}})$
rounding error, the accumulated absolute error is $O(N\eps_{\mathrm{mach}})$ without
compensation (or $O(\eps_{\mathrm{mach}})$ with Kahan summation~\cite{kahan1965}). We
switch to the Riemann--Siegel formula for $t > 10^3$
(Section~\ref{sec:rs}).

\begin{remark}[Multiplicity and Robustness of Detection]\label{rem:mult}
If a zeta zero $\rho_0$ has multiplicity $m > 1$ (widely conjectured not to occur, but
not proven), two effects must be distinguished:
\begin{enumerate}
  \item[(i)] The \emph{infinite series} $\eta(s)$ vanishes to order~$m$ at~$\rho_0$:
    $\eta(s) = c_m(s-\rho_0)^m + O(|s-\rho_0|^{m+1})$. This determines the \emph{width} of the dip in
    $|\eta(s)|$ around~$\rho_0$.
  \item[(ii)] The \emph{partial sum} $\eta_N(\rho_0) = O(N^{-\beta_0})$ by the Dirichlet
    test, \emph{independently of~$m$}.
\end{enumerate}
Consequently, the rate function jump is $\beta_0$, independent of~$m$. However, zeros 
of even multiplicity do not produce a sign change in $Z_H(t)$. To ensure a 
scientifically rigorous search that is robust against such cases (or GUE violations), 
our algorithm monitors local minima of the magnitude $|\mathcal{L}(\beta, t)|$ in 
addition to sign changes of $\mathcal{G}(\beta, t)$.
\end{remark}

\subsubsection{System~2: Generalized Loschmidt Amplitude}
\label{sec:system2}

The generalized Loschmidt amplitude is defined as
\begin{equation}\label{eq:G-def}
  \mathcal{G}(\beta,t) = \frac{1}{2Z(\beta)}
  \Bigl[e^{i\vartheta(t)}\,S_N(\beta+it)
  + e^{-i\vartheta(t)}\,S_N(\beta-it)\Bigr],
\end{equation}
where $\vartheta(t)$ is the Riemann--Siegel theta function:
\begin{equation}\label{eq:theta-def}
  \vartheta(t) = \im\!\bigl(\ln\Gamma(\tfrac{1}{4}+\tfrac{it}{2})\bigr)
  - \tfrac{t}{2}\ln\pi.
\end{equation}

Since $S_N(\beta-it) = \overline{S_N(\beta+it)}$ for real~$\beta$, we have
\begin{equation}\label{eq:G-real}
  \mathcal{G}(\beta,t) = \frac{1}{Z(\beta)}\,
  \re\!\bigl[e^{i\vartheta(t)}\,S_N(\beta+it)\bigr].
\end{equation}
While $Z_H(t)$ is exactly real in the limit $N \to \infty$ due to the functional 
equation, for finite~$N$, the partial sum $S_N$ does not satisfy the 
$\zeta(s) \leftrightarrow \zeta(1-s)$ symmetry. Consequently, 
$\mathcal{G}_N(\beta, t)$ is only an approximation whose deviations from 
the real axis (or the target function) scale as the tail of the Dirichlet series, 
typically $O(N^{-\beta})$. At $\beta=1/2$, this implies a remainder of 
$O(N^{-1/2})$, which must be accounted for in the validation tolerances.


\paragraph{Divergence of $Z(\beta)$ in System~2.}
The same divergence issue arises as in System~1. The resolution is identical: we work
with the finite-$N$ quantity $\mathcal{G}_N(\beta,t)$ and extract the rate function
$\mathcal{F}_2(\beta,t) = -\lim_{N\to\infty}(1/\ln N)\ln|\mathcal{G}_N(\beta,t)|$,
which has the same behavior as $\mathcal{F}_1$: singularities at zeta zeros, smooth
elsewhere.

\paragraph{Connection to the Hardy $Z$-function.}
At $\beta = \tfrac{1}{2}$, in the limit $N \to \infty$, the partial sum
$S_N(\tfrac{1}{2}+it)$ converges to $\zeta(\tfrac{1}{2}+it)$. The Riemann--Siegel
theta function is chosen precisely so that
\begin{equation}\label{eq:hardy-z}
  Z_H(t) = e^{i\vartheta(t)}\,\zeta\!\bigl(\tfrac{1}{2}+it\bigr)
\end{equation}
is real-valued for all real~$t$~\cite{hardy1914}. This is the Hardy $Z$-function. Its
sign changes correspond one-to-one with simple zeros of $\zeta(\tfrac{1}{2}+it)$ on the
critical line.

\paragraph{Why zeros of $\mathcal{G}$ correspond to zeta zeros.}
As with System~1, $Z_N(\tfrac{1}{2})$ diverges as $N \to \infty$, so
$\mathcal{G}_N(\tfrac{1}{2},t)$ itself tends to zero. The physically meaningful
statement concerns the \emph{ratio structure}: the numerator
$\re[e^{i\vartheta(t)}S_N(\tfrac{1}{2}+it)]$ converges to $Z_H(t)$ independently of
the divergent denominator. The zeros of $\mathcal{G}_N$ (detected at finite~$N$)
correspond to the zeros of $Z_H(t)$ in the limit. Since
$Z_H(t) = 0 \iff \zeta(\tfrac{1}{2}+it) = 0$ (because $e^{i\vartheta(t)} \ne 0$),
and Hardy~\cite{hardy1914} proved infinitely many zeros lie on the critical line,
System~2 detects exactly the nontrivial zeros on the critical line.

\begin{remark}[System~2 false positives at $\beta \ne \tfrac{1}{2}$]\label{rem:false-pos}
For $\beta \ne \tfrac{1}{2}$, $\mathcal{G}(\beta,t)$ can vanish even when
$\zeta(\beta+it) \ne 0$: since $e^{i\vartheta(t)}\zeta(\beta+it)$ is generally complex
for $\beta \ne \tfrac{1}{2}$, its real part vanishes whenever the phase of
$e^{i\vartheta(t)}\zeta(\beta+it)$ equals $\pm\pi/2$. These are \emph{phase-alignment
zeros}, not zeta zeros, and constitute false positives.

\emph{Consequence for experimental design:} System~2 is reliable for zero detection
\emph{only at $\beta = \tfrac{1}{2}$}. For off-critical-line scanning
($\beta \ne \tfrac{1}{2}$), System~1 (the $\mathcal{L}$ function, based on the eta
series) must be used. The experimental design (Section~\ref{sec:protocol}) is
structured accordingly.
\end{remark}


\subsection{The RH--DQPT Equivalence}
\label{sec:equivalence}

\begin{theorem}[RH--DQPT Equivalence]\label{thm:rh-dqpt}
The Riemann Hypothesis is equivalent to the following statement: for both System~1 and
System~2, dynamical quantum phase transitions (zeros of the order parameter /
non-analyticities of the rate function) occur only at inverse temperature
$\beta = \tfrac{1}{2}$.

More precisely:
\begin{enumerate}
  \item[(a)] RH holds if and only if, for all $t \in \R$ and all $\beta \in (0,1)$ with
    $\beta \ne \tfrac{1}{2}$: $\mathcal{L}(\beta,t) \ne 0$.
  \item[(b)] At $\beta = \tfrac{1}{2}$ specifically: RH holds if and only if the zeros of
    $\mathcal{G}(\tfrac{1}{2},t)$ (as a function of~$t$) account for \emph{all}
    nontrivial zeros of~$\zeta$. (System~2 cannot be used for off-line scanning due to
    false positives; statement~(a) via System~1 is the operative equivalence.)
\end{enumerate}
\end{theorem}

\begin{proof}
\emph{Part~(a).}
($\Rightarrow$) If RH holds, all nontrivial zeros have $\re(s) = \tfrac{1}{2}$. Then
$\eta(\beta+it) = (1-2^{1-(\beta+it)})\zeta(\beta+it) \ne 0$ for $\beta \ne \tfrac{1}{2}$
in $(0,1)$ (since $\zeta$ has no zeros there and the prefactor is nonzero for
$0 < \beta < 1$ as shown in Step~3). Hence $\mathcal{L}(\beta,t) \ne 0$.

($\Leftarrow$) If RH fails, there exists $\zeta(\beta_0+it_0) = 0$ with
$\beta_0 \ne \tfrac{1}{2}$ and $0 < \beta_0 < 1$. Then $\eta(\beta_0+it_0) = 0$, so
$\mathcal{L}(\beta_0,t_0) = 0$, giving a DQPT at $\beta_0 \ne \tfrac{1}{2}$.
Contrapositive: no DQPTs off the critical line implies~RH.
\end{proof}

\paragraph{What it would mean to find a DQPT at $\beta \ne \tfrac{1}{2}$.}
This would constitute a computational disproof of the Riemann Hypothesis: the existence
of a nontrivial zero of $\zeta(s)$ with $\re(s) \ne \tfrac{1}{2}$, with immediate
consequences for analytic number theory and the distribution of primes.


\subsection{Free Energy and Singularities}
\label{sec:free-energy}

The rate function (free energy density) for System~1 is
$\mathcal{F}_1(\beta,t) = -\lim_{N\to\infty}(1/\ln N)\ln|\mathcal{L}_N(\beta,t)|$,
as derived in Section~\ref{sec:system1}, Step~5. The values are:
\begin{align}
  \mathcal{F}_1(\beta,t) &= 1-\beta && \text{(away from zeros)}, \\
  \mathcal{F}_1(\beta,t_0) &= \beta_0+1-\beta && \text{(at a zero with $\re(s_0)=\beta_0$)}, \\
  \Delta\mathcal{F}_1 &= \beta_0 && \text{(DQPT jump magnitude)}.
\end{align}
At $\beta = \beta_0 = \tfrac{1}{2}$: $\mathcal{F}_1 = 1$ at the zero,
$\mathcal{F}_1 = \tfrac{1}{2}$ away. The DQPT jump is~$\tfrac{1}{2}$.

\begin{remark}[Normalization conventions]\label{rem:norm}
Wei et~al.~\cite{wei2025} use a different normalization employing $\log_2$ rather than
$\ln$. Under their convention, $\mathcal{F}_1^{\mathrm{Wei}} = \ln 2 \cdot
\mathcal{F}_1^{\mathrm{nat}}$, giving $\mathcal{F}_1^{\mathrm{Wei}} = \ln 2$ at zeros
and $(1-\beta)\ln 2$ away. We use natural logarithms throughout; numerical values differ
by a factor of~$\ln 2$.
\end{remark}

%% ====================================================================
\section{Numerical Methods}
\label{sec:methods}
%% ====================================================================

\subsection{Computation of Partial Dirichlet Series}
\label{sec:dirichlet}

The fundamental computational task is evaluating the partial Dirichlet series:
\begin{equation}\label{eq:partial-dirichlet}
  S_N(s) = \sum_{n=1}^{N} n^{-s}, \qquad s = \beta + it.
\end{equation}

\paragraph{Direct computation.}
Each term is computed as
\begin{equation}\label{eq:term-decomp}
  n^{-s} = n^{-\beta}\,n^{-it} = n^{-\beta}\,e^{-it\ln n},
\end{equation}
and separated into real and imaginary parts:
\begin{align}
  \re(n^{-s}) &= n^{-\beta}\cos(t\ln n), \label{eq:re-term}\\
  \im(n^{-s}) &= -n^{-\beta}\sin(t\ln n). \label{eq:im-term}
\end{align}
This decomposition avoids complex exponentiation and uses only real arithmetic (two
multiplications, one cosine, one sine per term), which is more numerically stable and
cache-friendly.

\paragraph{Alternating series.}
For System~1, we need the alternating partial sum
\begin{equation}\label{eq:eta-partial}
  \eta_N(s) = \sum_{n=1}^{N} (-1)^{n+1}\,n^{-s},
\end{equation}
obtained by multiplying each term by $(-1)^{n+1} = +1$ for odd~$n$, $-1$ for even~$n$.

\paragraph{Numerical stability.}
The terms $n^{-\beta}$ decrease monotonically for $\beta > 0$, ensuring numerical
stability of the running sum (no catastrophic cancellation in the absolute sense).
However, for the alternating series with small~$\beta$, consecutive terms nearly cancel,
leading to loss of relative precision. For $\beta = 0.5$ and large~$n$, consecutive
terms differ by $O(1/n)$, so the partial sum has relative precision degraded by a factor
of roughly $\sqrt{N}$. For $N = 10^5$, this costs approximately $2$--$3$ decimal digits,
which is acceptable in double precision ($\approx 15$--$16$ significant digits).

\paragraph{Precomputation.}
For a fixed~$\beta$, the values $n^{-\beta}$ ($n = 1, \ldots, N$) are independent of~$t$
and can be precomputed and stored. This reduces the per-point cost to $O(N)$
multiplications and trigonometric evaluations.

\paragraph{Error bound.}
The truncation error for the ordinary Dirichlet series is
\begin{equation}\label{eq:trunc-ordinary}
  |S_N(s) - \zeta(s)| \le \sum_{n=N+1}^{\infty} n^{-\beta}.
\end{equation}
By integral comparison:
\begin{equation}\label{eq:integral-bound}
  \sum_{n=N+1}^{\infty} n^{-\beta}
  \le \int_{N}^{\infty} x^{-\beta}\,dx
  = \frac{N^{1-\beta}}{\beta - 1}, \qquad \beta > 1.
\end{equation}
For $0 < \beta \le 1$ (the critical strip), the ordinary Dirichlet series does not
converge, but the alternating series converges. The truncation error for the alternating
series is bounded by the Dirichlet test (monotone coefficients, bounded partial sums of
$(-1)^{n+1}$):
\begin{equation}\label{eq:trunc-alt}
  |\eta_N(s) - \eta(s)| \le C(s)\,(N+1)^{-\beta},
\end{equation}
where $C(s) = O(1)$ for bounded $|s|$. For $\beta = 1/2$ and $N = 10^4$: error
$\lesssim 10^{-2}$; for $N = 10^6$: error $\lesssim 10^{-3}$. This is the fundamental
limitation of direct summation in the critical strip.

\paragraph{Euler--Maclaurin acceleration.}
For improved convergence, we apply the Euler--Maclaurin summation formula to the tail:
\begin{multline}\label{eq:euler-maclaurin}
  \sum_{n=N+1}^{M} n^{-s}
  \approx \int_{N}^{M} x^{-s}\,dx + \frac{1}{2}\bigl(N^{-s} + M^{-s}\bigr) \\
  + \sum_{k=1}^{p} \frac{B_{2k}}{(2k)!}\,
  s(s+1)\cdots(s+2k-2)\,\bigl(N^{-(s+2k-1)} - M^{-(s+2k-1)}\bigr) + R_p,
\end{multline}
where $B_{2k}$ are Bernoulli numbers and $R_p$ is the remainder. Taking $M \to \infty$
and using $p$~terms of the asymptotic expansion, the error bound must account for the
rising factorial $|s(s{+}1)\cdots(s{+}2p{-}2)|$, which grows as $|s|^{2p-1} \sim t^{2p-1}$
for large~$t$. The correct error estimate is
\begin{equation}\label{eq:em-error}
  |R_p| = O\bigl(|s|^{2p-1}\,N^{-(\beta+2p-1)}\bigr).
\end{equation}
This means that for large~$t$, the Euler--Maclaurin correction terms grow and the
asymptotic expansion eventually diverges (as is typical of asymptotic series). The
optimal number of correction terms $p$ depends on both $N$ and~$t$.

\paragraph{Precision table for Euler--Maclaurin correction ($\beta = 1/2$).}
In the following table, $|s|^{2p-1}$ denotes the magnitude of the rising factorial
$s(s{+}1)\cdots(s{+}2p{-}2)$, approximated as $|s|^{2p-1}$ for $|s| \gg 1$. With
$s = 1/2 + it$, we have $|s| \approx t$ for large~$t$.

\begin{table}[H]
\centering
\caption{Euler--Maclaurin effective error at $\beta = 1/2$.}
\label{tab:em-precision}
\begin{tabular}{@{}cccccc@{}}
\toprule
$N$ & $t$ & $p$ & $|s|^{2p-1}$ & $N^{-(\beta+2p-1)}$ & Effective error \\
\midrule
$10^4$ & $10^2$ & 3 & $10^{10}$  & $10^{-22}$ & $\sim 10^{-12}$ \\
$10^4$ & $10^2$ & 5 & $10^{18}$  & $10^{-38}$ & $\sim 10^{-20}$ \\
$10^4$ & $10^4$ & 3 & $10^{20}$  & $10^{-22}$ & $\sim 10^{-2}$  \\
$10^4$ & $10^4$ & 5 & $10^{36}$  & $10^{-38}$ & $\sim 10^{-2}$  \\
$10^5$ & $10^2$ & 5 & $10^{18}$  & $10^{-47.5}$ & $\sim 10^{-29.5}$ \\
$10^3$ & $10^2$ & 3 & $10^{10}$  & $10^{-16.5}$ & $\sim 10^{-6.5}$ \\
\bottomrule
\end{tabular}
\end{table}

The table reveals that for large~$t$, the Euler--Maclaurin correction becomes unreliable:
at $t = 10^4$ with any~$p$, the rising factorial overwhelms the $N^{-(\beta+2p-1)}$ decay.
This is the fundamental reason why the Riemann--Siegel formula (Section~\ref{sec:rs}) is
essential for $t > 10^3$. For moderate~$t$ ($t < 10^3$), Euler--Maclaurin with $p = 5$
and $N = 10^4$ provides $\sim 20$ digits of precision, far exceeding double-precision
limits.

\begin{remark}[On the relationship between acceleration and the DQPT framework]
\label{rem:acceleration-circularity}
A methodological subtlety deserves comment. The DQPT framework is formulated in terms of
the thermodynamic limit $N \to \infty$ of the \emph{finite} partial sums
$\mathcal{L}_N(\beta,t)$, and the physics resides in the finite-$N$ scaling behavior
(Section~\ref{sec:system1}). When we apply Euler--Maclaurin summation or Borwein
acceleration to compute the ``true'' value $\eta(s)$ (the infinite-sum limit), we are
effectively bypassing the finite-$N$ physics and returning to standard analytic number
theory. This does not invalidate the computation, but it means that the accelerated
results should be understood as \emph{high-precision evaluations of the same mathematical
function}, not as simulations of a physical system at finite~$N$. We use both approaches:
unaccelerated partial sums (to verify finite-$N$ scaling predictions of the DQPT theory)
and accelerated evaluations (to achieve the precision needed for reliable zero detection
and comparison with published tables).
\end{remark}

\paragraph{Borwein--Cohen--Villegas--Zagier acceleration.}
For the eta function specifically, Borwein~\cite{borwein2000} showed that the following
identity provides geometric convergence:
\begin{equation}\label{eq:borwein}
  \eta(s) = -\frac{1}{d_n}\sum_{k=0}^{n-1}(-1)^k\,(d_k - d_n)\,(k+1)^{-s},
\end{equation}
where the coefficients $d_k$ are defined by
\begin{equation}\label{eq:borwein-dk}
  d_k = n\sum_{j=0}^{k} \frac{(n+j-1)!\,4^j}{(n-j)!\,(2j)!}
\end{equation}
and $d_n = d_k|_{k=n}$.

\paragraph{Numerical stability of the coefficients.}
For $n = 40$, the coefficients $d_k$ grow extremely rapidly ($d_{40} \sim 10^{23}$),
exceeding the range of IEEE~754 double precision (which represents integers exactly only
up to $2^{53} \approx 9 \times 10^{15}$). To avoid overflow and catastrophic
cancellation, we compute the \emph{ratios} $c_k = (d_k - d_n)/d_n$ directly via a
stable recurrence. Define $e_k = d_k/d_n$; then $c_k = e_k - 1$, and the Borwein
formula becomes
\begin{equation}\label{eq:borwein-stable}
  \eta(s) = -\sum_{k=0}^{n-1}(-1)^k\,c_k\,(k+1)^{-s}.
\end{equation}
The ratios $e_k$ can be computed via the recurrence relation for the $d_k$ sequence,
working entirely in double precision without overflow (since $0 \le e_k \le 1$ for all
$k < n$).

\paragraph{Error bound and region of validity.}
The approximation error for the eta function satisfies~\cite{borwein2000}
\begin{equation}\label{eq:borwein-error}
  |\eta(s) - \text{approximation}| \le 3\,(3 + \sqrt{8}\,)^{-n}
\end{equation}
for $s$ in the region $\re(s) > 0$. With $n = 40$:
\begin{equation}\label{eq:borwein-40}
  |\text{error}_\eta| < 3 \times 5.828^{-40} \approx 10^{-30}.
\end{equation}
If one needs $\zeta(s)$ instead, it is recovered via $\zeta(s) = \eta(s)/(1 - 2^{1-s})$.
The resulting error in $\zeta$ is amplified by the factor $1/|1 - 2^{1-s}|$, which
diverges near $s = 1 + 2\pi ik/\log 2$ for $k \in \Z$. For computations in the critical
strip ($0 < \re(s) < 1$), $|1 - 2^{1-s}|$ is bounded away from zero. To see this:
$2^{1-s} = 1$ requires both $|2^{1-s}| = 2^{\re(1-s)} = 1$ (i.e., $\re(1-s) = 0$,
hence $\re(s) = 1$) and $\arg(2^{1-s}) = \im(1-s)\log 2 \in 2\pi\Z$. Since $\re(s) = 1$
lies outside the open critical strip $0 < \re(s) < 1$, the prefactor $1 - 2^{1-s}$ never
vanishes inside the strip. Moreover, for $\re(s) \in [0.1, 0.9]$,
$|1 - 2^{1-s}| \ge c > 0$ for some explicit constant~$c$ (computable for each~$\beta$),
so the error amplification is bounded.

This geometric convergence rate of $O(5.828^{-n})$ means $n = 40$ terms suffice for
$\sim 30$ digits of accuracy---far exceeding double precision requirements. This is our
primary method for computing $\eta(s)$ in the critical strip.

\paragraph{Kahan compensated summation.}
All summation loops use the Kahan compensated summation algorithm~\cite{kahan1965} to
reduce the accumulated floating-point rounding error from $O(N\,\eps_{\mathrm{mach}})$
to $O(\eps_{\mathrm{mach}})$, independent of~$N$. The algorithm maintains a running
compensation variable that tracks the low-order bits lost in each addition:
\begin{lstlisting}[language=C++,caption={Kahan compensated summation.}]
double sum = 0.0, c = 0.0;
for (int i = 0; i < N; ++i) {
    double y = x[i] - c;
    double t = sum + y;
    c = (t - sum) - y;
    sum = t;
}
\end{lstlisting}

\begin{remark}[FMA and GPU execution]\label{rem:fma}
CUDA compiles with fused multiply-add (FMA) instructions enabled by default. FMA
computes $a \cdot b + c$ with a single rounding (rather than the two roundings of
separate multiply and add), which changes the numerical result at the level of
$O(\eps)$ per operation. This has two consequences:
\begin{enumerate}
  \item The Kahan summation algorithm relies on the rounding behavior of individual
        additions. With FMA enabled, the compiler may fuse the compensation steps,
        potentially defeating the purpose of the algorithm. To ensure correct Kahan
        summation on GPU, use \texttt{\_\_dadd\_rn()} and \texttt{\_\_dsub\_rn()}
        intrinsics (round-to-nearest double-precision add/subtract) or compile with
        \texttt{--fmad=false} for the relevant kernels.
  \item CPU and GPU results will differ at $O(\eps)$ per operation due to different FMA
        behavior (CPU code compiled with \texttt{-ffp-contract=off} does not use FMA).
        When comparing CPU and GPU results for validation, the expected discrepancy is
        $O(N\,\eps)$ without Kahan summation, or $O(\eps)$ with correct Kahan summation
        on both platforms.
\end{enumerate}
\end{remark}


\subsection{Riemann--Siegel Theta Function}
\label{sec:theta}

The Riemann--Siegel theta function is defined by
\begin{equation}\label{eq:theta-def}
  \vartheta(t) = \im\!\bigl(\log\Gamma\bigl(\tfrac{1}{4} + \tfrac{it}{2}\bigr)\bigr)
  - \frac{t}{2}\log\pi.
\end{equation}

\paragraph{Stirling asymptotic expansion.}
For $|t| \gg 1$, we use the asymptotic expansion of $\log\Gamma$ via Stirling's series.
Setting $z = 1/4 + it/2$:
\begin{equation}\label{eq:stirling}
  \log\Gamma(z) = \Bigl(z - \tfrac{1}{2}\Bigr)\log z - z + \tfrac{1}{2}\log(2\pi)
  + \sum_{k=1}^{p} \frac{B_{2k}}{2k(2k-1)\,z^{2k-1}} + R_p(z).
\end{equation}
Taking the imaginary part and substituting $z = 1/4 + it/2$, we
obtain~\cite{edwards1974}\footnote{The first two correction terms ($1/(48t)$ and
$7/(5760t^3)$) are well known; the higher-order coefficients follow from the Stirling
series and are tabulated in Edwards~\cite{edwards1974}, Section~6.5.}
\begin{equation}\label{eq:theta-asymp}
  \vartheta(t) = \frac{t}{2}\ln\frac{t}{2\pi} - \frac{t}{2} - \frac{\pi}{8}
  + \frac{1}{48t} + \frac{7}{5760t^3} + \frac{31}{80640t^5}
  - \frac{127}{430080t^7} + O(t^{-9}).
\end{equation}
For our implementation, we use the first four correction terms (through $t^{-7}$), which
provides more than adequate precision for $t \ge 10$.

\paragraph{Error bound.}
The error after $p$ terms of the Stirling expansion is bounded by
\begin{equation}\label{eq:stirling-error}
  |R_p| \le \frac{|B_{2p+2}|}{(2p+2)(2p+1)\,|z|^{2p+1}}.
\end{equation}
For $t > 10$ and $p = 4$ (using terms through $t^{-7}$): $|R_4| < 10^{-12}$, which is
adequate for double precision. For $t > 50$ and $p = 3$ (through $t^{-5}$):
$|R_3| < 10^{-14}$.

\paragraph{Validity range.}
The asymptotic expansion is valid for $t > 0$ but achieves useful precision only for
$t > t_{\min}$. We require:
\begin{itemize}
  \item For double precision ($10^{-15}$): $t_{\min} \approx 10$ (with 4 correction terms).
  \item For extended precision ($10^{-30}$): $t_{\min} \approx 5$ (with 8 correction terms).
\end{itemize}
For $t < t_{\min}$, we fall back to numerical evaluation of $\log\Gamma$ using the
Lanczos approximation or a rational Chebyshev approximation for~$\Gamma$ in the relevant
region.

\begin{remark}[Branch cuts and small~$t$]\label{rem:branch}
The function $\vartheta(t) = \im(\log\Gamma(1/4 + it/2)) - (t/2)\log\pi$ involves the
complex logarithm of the Gamma function, which has branch cuts. For small~$t$ ($t < 10$),
the argument $1/4 + it/2$ is close to the real axis where $\Gamma$ has poles at
$0, -1, -2, \ldots$, and the principal branch of $\log\Gamma$ requires careful handling
to avoid discontinuities. Our implementation strategy:
\begin{itemize}
  \item For $t \ge 10$: use the Stirling asymptotic expansion ($4+$ correction terms),
        which is well away from branch cut issues and provides adequate precision.
  \item For $t < 10$: use the Lanczos approximation for $\Gamma(z)$ (with reflection
        formula for $\re(z) < 0.5$), or equivalently use the recurrence
        $\Gamma(z{+}1) = z\,\Gamma(z)$ to shift the argument into a region where the
        Lanczos series converges rapidly. This avoids the branch cut issues of
        $\log\Gamma$ by computing $\Gamma$ directly and then taking the logarithm.
\end{itemize}
In practice, the first zeta zero is at $t \approx 14.13$, so the small-$t$ regime is
only relevant for theta function evaluation at sub-zero-height positions (used in Gram
point computation and the Riemann--von Mangoldt formula).
\end{remark}


\subsubsection{The Riemann--Siegel Formula}
\label{sec:rs}

The Riemann--Siegel (RS) formula provides an efficient method for computing the Hardy
$Z$-function $Z_H(t) = e^{i\vartheta(t)}\zeta(1/2 + it)$ directly, requiring only
$O(\sqrt{t/(2\pi)})$ terms rather than $O(t)$ terms for direct
summation~\cite{titchmarsh1986,gabcke1979}. This is the standard method for large-scale
zero computations.

\paragraph{Explicit formula.}
The RS formula is
\begin{equation}\label{eq:rs-formula}
  Z_H(t) = 2\sum_{n=1}^{M} n^{-1/2}\cos\bigl(\vartheta(t) - t\ln n\bigr) + R(t),
\end{equation}
where $M = \lfloor\sqrt{t/(2\pi)}\rfloor$ is the main sum cutoff and $R(t)$ is the
remainder term.

\paragraph{First correction term.}
The leading-order remainder is given by the Riemann--Siegel
correction~\cite{gabcke1979}:
\begin{equation}\label{eq:rs-R0}
  R_0(t) = (-1)^{M-1}\Bigl(\frac{t}{2\pi}\Bigr)^{-1/4}\Psi(p),
\end{equation}
where $p = \sqrt{t/(2\pi)} - M$ is the fractional part and $\Psi(p)$ is the
Riemann--Siegel function:
\begin{equation}\label{eq:rs-psi}
  \Psi(p) = \frac{\cos\bigl(2\pi(p^2 - p - 1/16)\bigr)}{\cos(2\pi p)}.
\end{equation}
Higher-order correction terms $R_1, R_2, \ldots$ involve derivatives of~$\Psi$ and
powers of $(t/(2\pi))^{-1/2}$, each successive term reducing the error by an additional
factor of $O(t^{-1/2})$:
\begin{equation}\label{eq:rs-expansion}
  R(t) = R_0(t) + R_1(t) + R_2(t) + \cdots,
\end{equation}
where $|R_k(t)| = O(t^{-(2k+3)/4})$ (Siegel, 1932; Gabcke~\cite{gabcke1979}).

\paragraph{Error bounds after successive corrections.}
The following bounds are verified against
Titchmarsh~\cite{titchmarsh1986} and Gabcke~\cite{gabcke1979}:
\begin{align}
  \text{Main sum only:}\qquad &|Z_H(t) - \text{main sum}| = O(t^{-3/4}),
    \label{eq:rs-err0}\\
  \text{Main sum} + R_0:\qquad &|\text{error}| = O(t^{-5/4}),
    \label{eq:rs-err1}\\
  \text{Main sum} + R_0 + R_1:\qquad &|\text{error}| = O(t^{-7/4}).
    \label{eq:rs-err2}
\end{align}
Quantitatively (Gabcke~\cite{gabcke1979}): the error of the main sum alone is less than
$(3/2)\,t^{-3/4}$ for $t > 125$. Including the $R_0$ correction, the error is less than
$0.053\,t^{-5/4}$ for $t \ge 200$.

For $t = 10^4$: the uncorrected remainder is $O(t^{-3/4}) = O(10^{-3})$. Including
$R_0$, the error improves to $O(t^{-5/4}) = O(10^{-5})$. Including $R_0 + R_1$, the
error is $O(t^{-7/4}) = O(10^{-7})$.

\paragraph{Computational advantage.}
The main sum requires only $M = \lfloor\sqrt{t/(2\pi)}\rfloor$ terms:
\begin{center}
\begin{tabular}{@{}ll@{}}
\toprule
$t$ & $M$ \\
\midrule
$10^4$ & $\sim 40$ \\
$10^6$ & $\sim 399$ \\
$10^8$ & $\sim 3989$ \\
$10^{10}$ & $\sim 39894$ \\
\bottomrule
\end{tabular}
\end{center}
Compare this to direct summation of the Dirichlet series, which requires $N \gg t$ for
adequate convergence. The RS formula is essential for computations at large height.

\paragraph{Implementation notes.}
Each term in the main sum requires evaluation of
$n^{-1/2}\cos(\vartheta(t) - t\ln n)$. Since $n^{-1/2}$ is independent of~$t$, it can
be precomputed. The argument $\vartheta(t) - t\ln n$ must be reduced modulo $2\pi$
carefully to avoid loss of precision for large~$t$ (see Section~\ref{sec:trig} on
trigonometric evaluation).

\paragraph{Small-$t$ regime and method selection.}
The RS formula is an asymptotic approximation that improves with increasing~$t$.
At small heights, the main sum cutoff $M = \lfloor\sqrt{t/(2\pi)}\rfloor$ yields
very few terms:
\begin{center}
\begin{tabular}{@{}rrl@{}}
\toprule
$t$ & $M$ & RS viability \\
\midrule
$14$ & $1$ & unusable (single term) \\
$50$ & $2$ & poor \\
$100$ & $3$ & marginal \\
$200$ & $5$ & adequate (Gabcke bound applies) \\
$1000$ & $12$ & good \\
\bottomrule
\end{tabular}
\end{center}
For $t < 200$ (where the quantitative Gabcke error bound $0.053\,t^{-5/4}$ does not
apply), we compute the Hardy $Z$-function directly via
\begin{equation}\label{eq:zh-direct}
  Z_H(t) = \re\bigl[e^{i\vartheta(t)}\,\zeta\bigl(\tfrac{1}{2}+it\bigr)\bigr],
\end{equation}
where $\zeta(1/2+it)$ is obtained from the Borwein-accelerated eta function
(Section~\ref{sec:dirichlet}) as $\zeta(s) = \eta(s)/(1-2^{1-s})$. This provides
${\sim}30$ digits of precision independent of~$t$ (Paper Eq.~\ref{eq:borwein-40}).
For $t \ge 200$, the RS formula~\eqref{eq:rs-formula} with $R_0$ correction is used,
as it requires only $O(\sqrt{t})$ operations compared to the fixed cost of Borwein
($n = 100$ terms).

This two-method approach mirrors the strategy of Wei et~al.~\cite{wei2025}, who use
direct quantum simulation at small~$t$ and note that ``the deviation between estimated
and exact zeros decreases with increasing~$t$'' due to the improving RS approximation.
In our classical implementation, the Borwein method replaces the direct quantum
measurement for small~$t$.


\subsection{Zero Detection}
\label{sec:zero-detection}

\subsubsection{Bisection on the Hardy $Z$-function}
\label{sec:bisection}

The Hardy $Z$-function $Z_H(t)$ is real-valued, and its sign changes correspond to
simple zeros of $\zeta(1/2 + it)$. Our primary zero-detection method is:
\begin{enumerate}
  \item Evaluate $Z_H(t)$ on a uniform grid $t_0, t_0 + \Delta t, t_0 + 2\Delta t, \ldots$
  \item Identify intervals $[t_k, t_{k+1}]$ where $Z_H$ changes sign.
  \item Refine each sign change by bisection to precision~$\eps$.
\end{enumerate}

\paragraph{Grid spacing.}
The average spacing between consecutive zeros near height~$T$ is
\begin{equation}\label{eq:zero-spacing}
  \delta_{\mathrm{zero}} \sim \frac{2\pi}{\log(T/(2\pi))}.
\end{equation}
For $T = 100$: $\delta_{\mathrm{zero}} \approx 2.3$. For $T = 1000$:
$\delta_{\mathrm{zero}} \approx 1.2$. For $T = 10000$:
$\delta_{\mathrm{zero}} \approx 0.86$. We choose $\Delta t = \delta_{\mathrm{zero}}/4$
to ensure we do not miss any zeros (assuming zeros are roughly uniformly spaced, which
is a reasonable heuristic for the initial scan).

\subsubsection{Gram's Law and Its Failures}
\label{sec:gram}

Gram points are the values $g_n$ defined by $\vartheta(g_n) = n\pi$ for
$n = 0, 1, 2, \ldots$ Gram's ``law'' (actually a heuristic) states that
$(-1)^n Z_H(g_n) > 0$, i.e., $Z_H$ changes sign between consecutive Gram
points. This holds for approximately $73\%$ of Gram
intervals~\cite{trudgian2011}.

We use Gram points as an initial grid but do not rely on Gram's law for correctness.
Instead, we use finer subdivision wherever Gram's law fails (i.e., where $Z_H$ has the
same sign at two consecutive Gram points).

\subsubsection{Turing's Method for Zero Count Certification}
\label{sec:turing}

Turing~\cite{turing1953} developed a method to certify that all zeros up to height~$T$
have been found. The method relies on the Riemann--von Mangoldt formula for the
zero-counting function:
\begin{equation}\label{eq:rvm}
  N(T) = \frac{T}{2\pi}\log\frac{T}{2\pi} - \frac{T}{2\pi} + \frac{7}{8}
  + S(T) + O\!\left(\frac{1}{T}\right),
\end{equation}
where $S(T) = (1/\pi)\arg\zeta(1/2 + iT)$ satisfies $|S(T)| = O(\log T)$
unconditionally\footnote{See Trudgian~\cite{trudgian2011} for modern explicit bounds.
Note that $S(T)$ is \emph{not} bounded by a constant; the bound $|S(T)| = O(\log T)$
reflects the irregular distribution of zeta zeros.}. If the number of detected sign
changes of $Z_H$ in $[0, T]$ agrees with the nearest integer to
$(T/(2\pi))\log(T/(2\pi)) - T/(2\pi) + 7/8$, then all zeros in that range have been
found (up to the uncertainty introduced by $S(T)$).

We implement this check as a post-hoc validation: after detecting zeros up to
height~$T$, we verify that the count matches the Riemann--von Mangoldt prediction.

\subsubsection{Expected Zero Density}
\label{sec:zero-density}

The number of nontrivial zeros with $0 < \im(s) \le T$ is given by the Riemann--von
Mangoldt formula~\eqref{eq:rvm}. The leading terms give the approximation:
\begin{equation}\label{eq:rvm-approx}
  N(T) \approx \frac{T}{2\pi}\log\frac{T}{2\pi} - \frac{T}{2\pi} + \frac{7}{8}.
\end{equation}
This gives:
\begin{center}
\begin{tabular}{@{}rl@{}}
\toprule
$T$ & $N(T)$ \\
\midrule
$100$ & $\sim 29$ \\
$1{,}000$ & $\sim 649$ \\
$10{,}000$ & $\sim 10{,}142$ \\
$100{,}000$ & $\sim 138{,}069$ \\
\bottomrule
\end{tabular}
\end{center}


\subsection{Convergence Analysis}
\label{sec:convergence}

\subsubsection{Convergence of $L(\beta, t)$ with $N$}
\label{sec:conv-L}

For fixed $s = \beta + it$ with $0 < \beta < 1$, the alternating partial sum
$\eta_N(s)$ converges to $\eta(s)$ as $N \to \infty$. The rate of convergence is
\begin{equation}\label{eq:conv-rate}
  |\eta_N(s) - \eta(s)| = O(N^{-\beta}).
\end{equation}
This is the dominant source of error. For $\beta = 1/2$:
\begin{center}
\begin{tabular}{@{}rl@{}}
\toprule
$N$ & Error \\
\midrule
$10^4$ & $\sim 10^{-2}$ \\
$10^6$ & $\sim 10^{-3}$ \\
$10^8$ & $\sim 10^{-4}$ \\
\bottomrule
\end{tabular}
\end{center}
This slow convergence is the fundamental computational bottleneck. To achieve precision
$\delta$, we need $N \sim \delta^{-1/\beta} = \delta^{-2}$ for $\beta = 1/2$. This is
much worse than the Riemann--Siegel formula, which achieves precision~$\delta$ with
$O(\sqrt{T})$ terms, but the direct approach has the virtue of simplicity and direct
connection to the physical (DQPT) interpretation.

\paragraph{Improvement via Euler--Maclaurin.}
With $p$ terms of Euler--Maclaurin correction applied to the tail sum beyond~$N$, the
effective error becomes $O(|s|^{2p-1}\,N^{-(\beta+2p-1)})$ (see
Section~\ref{sec:dirichlet} for the full bound and precision table). For $p = 5$,
$\beta = 1/2$, and moderate~$t$ ($t < 10^4$): the error is well below double-precision
limits for $N \ge 10^4$. This makes Euler--Maclaurin-corrected summation practical for
moderate~$t$. For large~$t$, the Riemann--Siegel formula (Section~\ref{sec:rs}) is
preferable.

\subsubsection{Required $N$ as a Function of $t$ and Precision}
\label{sec:required-N}

The terms $n^{-s} = n^{-\beta}e^{-it\ln n}$ oscillate with frequency proportional
to~$t$. The Nyquist criterion requires that we sample these oscillations adequately, but
since we are summing (not sampling) the series, the oscillation frequency does not
directly increase the required~$N$. However, near zeros of~$\zeta$, the partial sums
exhibit cancellation at scale $O(N^{-\beta})$, so detecting zeros to relative
precision~$\delta$ requires
\begin{equation}\label{eq:required-N}
  N \ge \delta^{-1/\beta}.
\end{equation}
For $\beta = 1/2$ and $\delta = 10^{-6}$ (sufficient to distinguish a zero from a
near-miss): $N \ge 10^{12}$. This is impractical without acceleration techniques.

With Euler--Maclaurin ($p = 5$) and accounting for the $|s|^{2p-1}$ factor: the required
$N$ satisfies $|s|^{2p-1}\,N^{-(\beta+2p-1)} < \delta$, giving
\begin{equation}\label{eq:required-N-em}
  N \ge \Bigl(\frac{|s|^{2p-1}}{\delta}\Bigr)^{1/(\beta+2p-1)}.
\end{equation}
For $t = 100$, $|s| \approx 100$, $p = 5$, $\beta = 0.5$:
$N \ge (10^{18}/10^{-6})^{1/9.5} = 10^{24/9.5} \approx 336$. For $t = 10^4$:
$N \ge (10^{36}/10^{-6})^{1/9.5} = 10^{42/9.5} \approx 26{,}300$. These are achievable,
showing that Euler--Maclaurin correction is essential for practical computation at high
precision, though the required $N$ grows with~$t$.

\subsubsection{Computational Complexity}
\label{sec:complexity}

Direct summation: $O(N)$ operations per $(\beta, t)$ point. For $M$ grid points in the
$(\beta, t)$ plane, the total cost is $O(N \cdot M)$.

With Euler--Maclaurin: $O(N + p)$ per point, where $p$ is small ($5$--$10$). For a
moderate $N$ ($10^3$ to $10^4$), the cost is dominated by the direct sum, giving
$O(N \cdot M)$ total.

For a grid of $M_\beta = 100$ values of~$\beta$ and $M_t = 10{,}000$ values of~$t$,
with $N = 10^4$:
\begin{equation}\label{eq:total-ops}
  \text{Total operations} \sim 10^4 \times 10^6 = 10^{10}.
\end{equation}
On a single CPU core, this workload requires substantial time (minutes to hours 
depending on architecture and instruction set). While GPU parallelism 
drastically reduces wall-clock time, the hardware-level bottleneck of 
FP64 transcendental units must be considered (Section~\ref{sec:throughput}).


\subsection{GPU Parallelization Strategy}
\label{sec:gpu}

\subsubsection{Parallelism Structure}
\label{sec:parallel}

The computation of $L(\beta, t)$ or $G(\beta, t)$ at each $(\beta, t)$ point is
completely independent of all other points. This is an embarrassingly parallel workload,
ideal for GPU execution.

\paragraph{Parallelization scheme.}
\begin{itemize}
  \item Each CUDA thread computes one $(\beta, t)$ point.
  \item The thread iterates over $n = 1, \ldots, N$, accumulating the partial sum.
  \item The result (a complex number for~$L$, a real number for~$G$) is written to global
        memory.
\end{itemize}

\subsubsection{Memory Model}
\label{sec:memory}

For each $\beta$ value, we precompute the array $\{n^{-\beta}\}_{n=1}^{N}$ and store it
in global memory (or texture memory for cached access). Each thread reads this array and
computes the trigonometric factors $\cos(t\ln n)$ and $\sin(t\ln n)$ on the fly.

\paragraph{Memory requirements.}
\begin{itemize}
  \item Precomputed $n^{-\beta}$: $N \times \texttt{sizeof(double)} = N \times 8$ bytes
        per $\beta$ value. For $N = 10^5$ and $100$~$\beta$ values: $80$~MB (fits in
        $6$~GB VRAM).
  \item Precomputed $\ln(n)$: $N \times 8$ bytes $= 0.8$~MB (shared across all $\beta$
        values).
  \item Output array: $M_\beta \times M_t \times \texttt{sizeof(double)}$ per output
        quantity. For $100 \times 10{,}000 = 10^6$ points: $8$~MB.
  \item Total VRAM usage: $< 100$~MB, well within $6$~GB.
\end{itemize}

\subsubsection{Expected Throughput}
\label{sec:throughput}

RTX~2060 specifications:
\begin{itemize}
  \item 1920 CUDA cores.
  \item 6~GB GDDR6 VRAM.
  \item Base clock: 1365~MHz.
  \item FP64 throughput: approximately $210$~GFLOPS ($1/32$ of FP32).
\end{itemize}

Since our computation involves double-precision transcendental functions (\texttt{cos},
\texttt{sin}, \texttt{exp}, \texttt{log}), the throughput is limited by the special
function units (SFU).

\begin{remark}[SFU throughput for transcendentals]\label{rem:sfu}
On NVIDIA GPUs, the SFU throughput for double-precision transcendental functions is
approximately $1/4$ of the FP64 ALU throughput (and for some transcendentals, up to
$1/10$ due to multi-pass evaluation). Since our workload is transcendental-heavy (each
term requires one \texttt{cos} and one \texttt{sin} evaluation), the effective throughput
must be derated by $4$--$10\times$ compared to the peak FP64 rate.
\end{remark}

\paragraph{Revised throughput estimates.}
\begin{itemize}
  \item Per-point cost: $N$ multiplications $+ N$ transcendental evaluations. The
        transcendentals dominate: effective cost $\sim 10N$ FLOPS equivalent (accounting
        for SFU pipeline occupancy).
  \item For $N = 10^4$: $\sim 10^5$ FLOPS-equivalent per point.
  \item Effective throughput:
        $(210 \times 10^9/4)/(10^5) \approx 5 \times 10^5$ points/second (conservative;
        actual throughput depends on occupancy and memory access patterns).
\end{itemize}

For a $180 \times 10{,}000$ grid ($1.8 \times 10^6$ points): approximately $3$--$4$
seconds. For a $180 \times 100{,}000$ grid ($1.8 \times 10^7$ points): approximately
$30$--$40$ seconds. For $N = 10^5$ and a $180 \times 10{,}000$ grid: approximately
$30$--$40$ seconds.

These revised estimates ($4$--$10\times$ slower than naive FP64 peak estimates) still show
that GPU-accelerated computation is more than sufficient for our parameter ranges. Even
with additional overheads (memory latency, register pressure, thread divergence), runtimes
remain under a few minutes for the full sweep.

\subsubsection{Batch Size Estimation}
\label{sec:batch}

We organize the computation into batches to manage VRAM and kernel launch overhead:
\begin{itemize}
  \item Batch size: $10^5$ to $10^6$ $(\beta, t)$ points per kernel launch.
  \item Number of batches: total grid size / batch size.
  \item For a $100 \times 100{,}000$ grid: $10^7 / 10^6 = 10$ batches.
\end{itemize}
Each batch completes in $< 1$~second, so the total wall-clock time (including data
transfer) is under $10$~seconds for the full sweep.

%% ====================================================================
\section{Experimental Design}
\label{sec:experiment}
%% ====================================================================

\subsection{Hypothesis}
\label{sec:hypothesis}

We formulate the investigation as a deterministic numerical verification with the
following decision framework.\footnote{Unlike a classical Neyman--Pearson test, there is
no randomness here---the computation is deterministic. We adopt the hypothesis language
for clarity of exposition, not to claim statistical inference.}

\paragraph{$H_0$ (Null hypothesis).}
The Riemann Hypothesis holds. All DQPTs in both System~1 and System~2 occur exclusively
at $\beta = 1/2$. Equivalently, for all $t \in \R$ and all $\beta \in (0,1)$ with
$\beta \neq 1/2$: $L(\beta, t) \neq 0$ and $G(\beta, t)$ does not exhibit a zero
corresponding to a zeta zero.

\paragraph{$H_1$ (Alternative hypothesis).}
The Riemann Hypothesis is false. There exists at least one DQPT at some
$\beta_0 \neq 1/2$, corresponding to a nontrivial zero of $\zeta(s)$ with
$\re(s) = \beta_0 \neq 1/2$.

We note that failure to reject $H_0$ does not constitute proof of~RH; it merely provides
additional numerical evidence consistent with~RH up to the scanned range.


\subsection{Test Protocol}
\label{sec:protocol}

The investigation proceeds in four stages.

\paragraph{Stage~1: Implementation Verification.}
Compute $L(1/2, t_k)$ and $G(1/2, t_k)$ at the first $100$ known nontrivial zeta zeros
($t_1 = 14.134725\ldots$, $t_2 = 21.022040\ldots$, etc.) and verify that they vanish
within numerical tolerance. This confirms that the implementation correctly reproduces
the theoretical predictions.

\paragraph{Stage~2: Critical-Line Zero Detection.}
Scan $t \in [0, T]$ at fixed $\beta = 1/2$, detecting zeros of $G(1/2, t)$ by sign
changes. Compare the detected zero positions with published
tables~\cite{odlyzko1989,lmfdb}. Verify that the zero count matches the Riemann--von
Mangoldt prediction $N(T)$.

\paragraph{Stage~3: Off-Critical-Line Scan (System~1 only).}
For each $\beta \in [0.1, 0.9]$ with spacing $\Delta\beta = 0.005$, excluding
$[0.4975, 0.5025]$, scan $t \in [0, T]$ and compute $|L(\beta, t)|$ using System~1 (the
eta function). System~2 (the $G$~function) is \emph{not} used for off-line scanning
because it can produce false positives at $\beta \neq 1/2$ due to phase-alignment zeros
(see Section~\ref{sec:system2}). Record the minimum value of $|L(\beta, t)|$ over the
$t$-range for each~$\beta$. If this minimum is below the zero-detection threshold~$\eps$,
flag it for further investigation (Stage~4). Apply the adaptive refinement strategy
described in Section~\ref{sec:params} (Grid resolution).

\paragraph{Stage~4: High-Precision Refinement.}
For any candidate $(\beta, t)$ flagged in Stage~3 with $|L(\beta, t)| < \eps$:
\begin{enumerate}
  \item[(a)] Increase $N$ and recompute to distinguish a genuine zero from a near-miss.
  \item[(b)] Refine $(\beta, t)$ by 2D Newton's method or simplex optimization.
  \item[(c)] If the refined $\beta$ is within tolerance of $1/2$, classify as a
             critical-line zero (consistent with~$H_0$). If $\beta$ remains bounded away
             from $1/2$, classify as a potential counterexample to~RH (rejecting~$H_0$).
\end{enumerate}


\subsection{Parameters}
\label{sec:params}

\paragraph{Truncation parameter $N$.}
\begin{itemize}
  \item \emph{Phase~1 (verification):} $N = 10^4$ with Euler--Maclaurin correction
        ($5$~terms). Justification: effective precision $\sim 10^{-38}$, far exceeding
        double-precision limits.
  \item \emph{Phase~3 (scanning):} $N = 10^3$ to $10^4$, depending on available compute.
        Justification: for a coarse scan, precision $\sim 10^{-2}$ suffices to detect
        zeros; flagged candidates are refined in Phase~4.
  \item \emph{Phase~4 (refinement):} $N = 10^5$ with Euler--Maclaurin correction.
\end{itemize}

\paragraph{Grid resolution.}
\begin{itemize}
  \item $\Delta\beta = 0.005$ for the coarse scan (at least $179$ off-critical-line
        $\beta$ values in $[0.1, 0.9]$ excluding a neighborhood of~$0.5$). The finer
        grid (compared to the naive $\Delta\beta = 0.05$) is necessary because a zero at
        $\beta_0 = 0.52$ would produce $|\eta(0.50 + it_0)| \sim 0.02$, which is too
        small to reliably detect at a grid point $0.02$ away. With $\Delta\beta = 0.005$,
        such a zero would be sampled at $\beta = 0.52$ (or within $0.0025$ of it), where
        $|\eta|$ is detectably small.
  \item $\Delta\beta = 0.001$ for refined scans near any flagged candidate.
  \item $\Delta t = 0.1$ for $t \in [0, 1000]$ ($10^4$ points per $\beta$ value).
  \item $\Delta t = 0.01$ for refined scans near flagged candidates.
\end{itemize}

\paragraph{Adaptive refinement strategy.}
Zeros with $|\beta_0 - 1/2| < \delta_{\mathrm{threshold}}$ may produce only marginal
dips in $|L(\beta, t)|$ at nearby grid points. To address this:
\begin{enumerate}
  \item \emph{Initial coarse scan:} Evaluate $|L(\beta, t)|$ on the
        $\Delta\beta = 0.005$ grid and record all local minima below a liberal threshold
        ($\eps_{\mathrm{coarse}} = 10^{-3}$).
  \item \emph{Adaptive refinement:} For each flagged $(\beta, t)$ region, bisect in both
        $\beta$ and~$t$ to find the true minimum of $|L|$. If the minimum decreases with
        refinement (indicating a genuine zero rather than a near-miss), continue refining
        until $|L| < 10^{-6}$ or the minimum stabilizes.
  \item \emph{Convergence check:} Increase $N$ and verify that $|L|$ continues to
        decrease (as expected for a genuine zero) rather than stabilizing (indicating a
        near-miss).
\end{enumerate}

\paragraph{Zero-detection tolerance.}
\begin{itemize}
  \item $\eps = 10^{-6}$ for classifying a point as a zero (coarse scan).
  \item $\eps = 10^{-12}$ for refined confirmation.
\end{itemize}

\paragraph{Scan range.}
\begin{itemize}
  \item $T_{\max} = 1000$ for the initial investigation (capturing $\sim 649$ zeros on
        the critical line).
  \item $T_{\max} = 10{,}000$ for extended runs ($\sim 10{,}142$ zeros).
  \item $T_{\max} = 100{,}000$ as an aspirational target with GPU acceleration
        ($\sim 138{,}069$ zeros).
\end{itemize}

\paragraph{Hardware.}
\begin{itemize}
  \item CPU: specified at runtime; any modern x86-64 processor suffices.
  \item GPU: NVIDIA RTX~2060 (1920 CUDA cores, 6~GB GDDR6 VRAM, Compute
        Capability~7.5).
  \item Precision: IEEE~754 double precision (64-bit, $\sim 15.9$ significant decimal
        digits).
\end{itemize}


\subsection{Validation}
\label{sec:validation}

\subsubsection{Unit Tests Against Reference Values}
\label{sec:unit-tests}

Each core mathematical function is tested against high-precision reference values
computed independently using \texttt{mpmath} (Python arbitrary-precision library) or
Mathematica. Specifically:
\begin{itemize}
  \item $S_N(s)$ for $N = 100$ and $s = 0.5 + 14.134725i$: compare to \texttt{mpmath}.
  \item $\vartheta(t)$ for $t = 14.134725$, $21.022040$, $100$, $1000$: compare to
        \texttt{mpmath}.
  \item $\eta(s)$ for $s = 0.5 + 14.134725i$: compare to \texttt{mpmath}.
  \item $Z_H(t)$ for $t = 14.134725$ (first zero): should be zero within precision.
\end{itemize}
Tolerance: agreement to $10$ significant digits (limited by double precision and~$N$).

\subsubsection{Cross-Validation: System~1 vs System~2}
\label{sec:cross-val}

Both systems must yield the same zero positions (to within numerical precision) at
$\beta = 1/2$. Specifically, if $L(1/2, t^*) = 0$, then $G(1/2, t^*)$ must also change
sign near $t = t^*$. We verify this for all detected zeros.

Tolerance: $|t^*_L - t^*_G| < 10^{-4}$ for each corresponding zero pair.

\subsubsection{Convergence Tests}
\label{sec:conv-tests}

For each of several fixed $(\beta, t)$ points, compute $L(\beta, t)$ for
$N = 10^2, 10^3, 10^4, 10^5$ and verify that the values converge. Specifically:
\begin{equation}\label{eq:conv-test}
  |L_N(\beta, t) - L_{N/10}(\beta, t)| < C \cdot (N/10)^{-\beta},
\end{equation}
where $C$ is a constant. The convergence rate should be consistent with the theoretical
$O(N^{-\beta})$.

\subsubsection{Comparison with Odlyzko's Tables}
\label{sec:odlyzko-compare}

The first $100$ detected zeros are compared with Odlyzko's published tables (available
from the LMFDB~\cite{lmfdb} and Andrew Odlyzko's
website~\cite{odlyzko1989,odlyzko2001}). Agreement is required to within the precision
of our computation.


\subsection{Controls}
\label{sec:controls}

The following consistency checks serve as controls to verify the correctness of the
implementation.

\paragraph{Control~1.}
At known zeros ($t = t_k$), verify:
\begin{equation}\label{eq:control1}
  F_1(1/2, t_k) = 1 \pm \eps,
\end{equation}
where $\eps = 0.05$ (accounting for finite-$N$ effects).\footnote{Under the Wei et~al.\
$\log$-base-$2$ normalization, this value would be $\ln 2 \approx 0.693$; we use the
natural log normalization throughout this document, giving
$F_1 = \beta_0 + 1 - \beta = 1$ at $\beta = \beta_0 = 1/2$.}

\paragraph{Control~2.}
Away from zeros, for generic~$t$ (e.g., $t = 50.0$, which is not near any zero), verify:
\begin{equation}\label{eq:control2}
  F_1(\beta, t) = (1 - \beta) \pm \eps
\end{equation}
for $\beta \in \{0.3, 0.5, 0.7\}$.\footnote{Under Wei et~al.\ normalization:
$(1 - \beta)\ln 2$.}

\paragraph{Control~3.}
Verify the zero count:
\begin{equation}\label{eq:control3}
  |N_{\mathrm{detected}}(T) - N_{\mathrm{predicted}}(T)| \le 1,
\end{equation}
where $N_{\mathrm{predicted}}(T) = \mathrm{round}((1/\pi)\vartheta(T) + 1)$. The
tolerance of~$1$ accounts for the $S(T) = (1/\pi)\arg\zeta(1/2 + iT)$ fluctuation term,
which can cause the actual zero count to differ from the smooth prediction by $\pm 1$ at
specific $T$ values.

\paragraph{Control~4.}
Verify that $G(1/2, t)$ is real-valued (imaginary part $< 10^{-14}$) for all sampled $t$
values, confirming the $Z$-function property.

\paragraph{Control~5.}
Verify the functional equation numerically: for selected $s$ values, confirm that
$\zeta(s)$ and $\zeta(1-s)$ satisfy the functional equation~\eqref{eq:func-eq} within
precision. (This tests the Riemann--Siegel theta function implementation indirectly.)

%% ====================================================================
\section{Error Analysis}
\label{sec:errors}
%% ====================================================================

\subsection{Truncation Error (Finite $N$)}
\label{sec:trunc}

The dominant systematic error arises from truncating the Dirichlet series at $N$ terms.

\paragraph{Without Euler--Maclaurin correction.}
\begin{equation}\label{eq:trunc-no-em}
  |\eta_N(s) - \eta(s)| \le C(s)\,(N{+}1)^{-\beta},
\end{equation}
where $C(s) = O(1)$ for bounded $|s|$. For $\beta = 0.5$ and $N = 10^4$: truncation
error $\sim 10^{-2}$.

\paragraph{With Euler--Maclaurin correction ($p$ terms).}
\begin{equation}\label{eq:trunc-em}
  |\eta_N^{\mathrm{EM}}(s) - \eta(s)|
  \le C_p\,|s|^{2p-1}\,N^{-(\beta + 2p - 1)},
\end{equation}
where the factor $|s|^{2p-1}$ accounts for the rising factorial in the correction terms
(see Section~\ref{sec:dirichlet}). For $\beta = 0.5$, $p = 5$, $N = 10^4$, and
$t = 100$: truncation error $\sim 10^{9} \times 10^{-38} = 10^{-29}$. See the precision
table in Section~\ref{sec:dirichlet} (Table~\ref{tab:em-precision}) for additional
parameter combinations.

This error is purely systematic and can be reduced by increasing $N$ (at the cost of
more computation) or by choosing $p$ optimally for the given $(N, t)$ pair.


\subsection{Floating-Point Precision}
\label{sec:trig}

IEEE~754 double-precision arithmetic provides approximately $15.9$ significant decimal
digits ($53$-bit mantissa). The key sources of floating-point error are:

\paragraph{Summation error.}
Accumulating $N$ terms in floating point introduces rounding error of order
$N\,\eps_{\mathrm{mach}} \sim N \times 10^{-16}$. For $N = 10^4$: rounding error
$\sim 10^{-12}$. For $N = 10^6$: rounding error $\sim 10^{-10}$. This can be mitigated
by Kahan compensated summation~\cite{kahan1965}, which reduces the error to
$O(\eps_{\mathrm{mach}})$ independent of~$N$.

\paragraph{Trigonometric evaluation and the FP64 precision wall.}
$\cos(t\ln n)$ and $\sin(t\ln n)$ are evaluated by the standard library, with error
$O(\eps_{\mathrm{mach}})$ per evaluation. For large arguments ($t\ln n \gg 1$), the
argument reduction step introduces additional error proportional to
$t\ln(n)\,\eps_{\mathrm{mach}}$. This constitutes a fundamental precision wall for
double-precision computation:
\begin{center}
\begin{tabular}{@{}rll@{}}
\toprule
$t$ & Argument $t\ln n$ (at $n = 10^4$) & Effective digits lost \\
\midrule
$10^3$ & $\sim 10^4$ & $\sim 4$ \\
$10^5$ & $\sim 10^6$ & $\sim 6$ \\
$10^7$ & $\sim 10^8$ & $\sim 8$ \\
$10^{10}$ & $\sim 10^{11}$ & $\sim 11$ \\
\bottomrule
\end{tabular}
\end{center}
At $t > 10^7$, fewer than $8$ significant digits remain in the trigonometric evaluation,
degrading zero detection to tolerances worse than $10^{-8}$. At $t > 10^{10}$, the
argument reduction error ($\sim 10^{-5}$) exceeds the zero-detection threshold, rendering
double-precision computation unreliable.

For scientifically rigorous computation at large heights ($t > 10^7$), arbitrary-precision
arithmetic (e.g., MPFR~/ GNU~MP, or at minimum IEEE~754 quad precision via
\texttt{\_\_float128} or \texttt{long double} on x86) is required. The Riemann--Siegel
formula (Section~\ref{sec:rs}) partially mitigates this issue by requiring only
$M = O(\sqrt{t})$ terms (reducing the maximum argument to
$t\ln M \sim t\ln\sqrt{t} = (t/2)\ln t$), but the theta function evaluation
$\vartheta(t)$ itself involves arguments of order~$t$, so the precision wall persists.
Our current implementation is therefore limited to $t < 10^7$ for reliable zero detection
in double precision, with reduced reliability up to $t \sim 10^{10}$.

\paragraph{Cancellation near zeros.}
When $L(\beta, t)$ is near zero, the partial sum involves extensive cancellation. The
absolute value of~$L$ near a zero is $O(N^{-\beta})$, while individual terms have
magnitude $O(n^{-\beta})$ with the largest terms $O(1)$. Without compensated summation,
the absolute error in the sum is $O(N\,\eps_{\mathrm{mach}})$ (each of $N$ additions
contributes $O(\eps_{\mathrm{mach}})$ rounding error). With Kahan summation, the absolute
error is reduced to $O(\eps_{\mathrm{mach}})$, independent of~$N$. The value of~$L$ near
a zero is $O(N^{-\beta})$. The relative precision is therefore:
\begin{align}
  \text{Without Kahan:}\quad &\text{relative error}
    \sim \frac{N\,\eps_{\mathrm{mach}}}{N^{-\beta}}
    = N^{1+\beta}\,\eps_{\mathrm{mach}}, \label{eq:no-kahan}\\
  \text{With Kahan:}\quad &\text{relative error}
    \sim \frac{\eps_{\mathrm{mach}}}{N^{-\beta}}
    = N^{\beta}\,\eps_{\mathrm{mach}}. \label{eq:with-kahan}
\end{align}
For $N = 10^4$ and $\beta = 0.5$: without Kahan, relative precision
$\sim 10^6 \times 10^{-16} = 10^{-10}$. With Kahan summation:
$\sim N^\beta\,\eps_{\mathrm{mach}} = 10^2 \times 10^{-16} = 10^{-14}$.

This is adequate for zero detection with tolerance $10^{-6}$.


\subsection{Underflow Guard for $\log|L_N|$}
\label{sec:underflow}

The rate function $F_1(\beta, t) = -(1/\ln N)\ln|L_N(\beta, t)|$ involves taking the
logarithm of $|L_N|$. In floating-point arithmetic, $|L_N|$ can underflow to exactly
zero (when $\beta$ is large or at a near-zero of $\eta$ with large~$N$), producing
$\log(0) = -\infty$ and corrupting downstream computations.

\paragraph{Guard implementation.}
We use:
\begin{equation}\label{eq:underflow-guard}
  F_1 = -\frac{1}{\ln N}\,\ln\bigl(\max\bigl(|L_N|,\, 10^{-300}\bigr)\bigr).
\end{equation}
The floor value $10^{-300}$ is well above the IEEE~754 double-precision minimum subnormal
($\sim 5 \times 10^{-324}$) but far below any physically meaningful value. When
$|L_N| < 10^{-300}$, the computed $F_1$ is flagged as a sentinel value indicating
numerical underflow, and the point is marked for investigation with higher precision or
larger~$N$.

Additionally, any $(\beta, t)$ point where $|L_N| = 0$ exactly (to double precision) is
recorded separately and not included in the rate function analysis, since it may
represent either a genuine zero (signaling a DQPT) or a floating-point artifact.


\subsection{Grid Discretization Error}
\label{sec:grid-error}

Zeros detected by bisection on a grid with spacing $\Delta t$ have position error bounded
by $\Delta t/2$ in the initial detection, refined to arbitrary precision by iterative
bisection. After $k$ bisection iterations, the position error is
$\Delta t \cdot 2^{-(k+1)}$. For $\Delta t = 0.1$ and $k = 30$ iterations: position
error $\sim 10^{-10}$.

In the $\beta$ direction, the grid spacing $\Delta\beta$ determines our ability to
localize off-critical-line zeros. A zero at $\beta_0$ will be detected (in the sense
that $|L|$ dips below~$\eps$) only if $\beta_0$ is within $\sim\!\Delta\beta$ of a grid
point. For $\Delta\beta = 0.05$, we could miss a zero with $\beta_0 = 0.52$ (only $0.02$
from the critical line). However, the minimum of $|L(\beta, t)|$ over~$t$ is a smooth
function of~$\beta$, so a zero at $\beta_0 = 0.52$ would also produce a small value at
the nearest grid point $\beta = 0.50$ (the critical line)---making it detectable, though
its $\beta$-position would initially be misattributed to~$0.50$. The refinement stage
(Stage~4) would then resolve its true position.


\subsection{False Positive and False Negative Rates}
\label{sec:false-rates}

\paragraph{False positives} (declaring a zero where none exists) arise from numerical
artifacts such as:
\begin{itemize}
  \item Two nearly-cancelling large terms producing a spuriously small sum.
  \item Floating-point underflow near a local minimum of $|L|$.
\end{itemize}
Mitigation: every detected zero is confirmed by (a)~increasing~$N$, (b)~computing both
$L$ and~$G$, and (c)~verifying consistency with the Riemann--von Mangoldt zero count.

\paragraph{False negatives} (missing a genuine zero) arise from:
\begin{itemize}
  \item Grid spacing too coarse relative to the zero separation.
  \item Insufficient $N$ causing the partial sum to not resolve the zero.
\end{itemize}
Mitigation: grid spacing is chosen to be at most $1/4$ of the expected zero separation
(Section~\ref{sec:bisection}). The zero count check (Section~\ref{sec:controls},
Control~3) provides an independent confirmation that no zeros were missed.

%% ====================================================================
\section{Expected Results}
\label{sec:results}
%% ====================================================================

\subsection{If RH Holds (Expected Scenario)}
\label{sec:rh-holds}

Under the assumption that~RH is true---supported by all prior numerical evidence and by
classical zero-density estimates, which imply that $100\%$ of zeros lie on the critical
line in a density sense\footnote{This density result already followed from the prior
$O(T^{3/4+\eps})$ bound on the number of off-line zeros~\cite{ingham1940}, since
$T^{3/4}/N(T) \to 0$. Guth and Maynard's 2024 improvement~\cite{guthmaynard2024}
sharpens the exponent to $O(T^{3/5+\eps})$, the first improvement in over 80~years, but
the qualitative conclusion of $100\%$ density was already known.}---we expect:

\begin{enumerate}
  \item \textbf{System~1 ($L$):} $L(1/2, t)$ vanishes at each of the $\sim 649$
        nontrivial zeros with $0 < \im(s) \le 1000$, with $|L(1/2, t_k)| < \eps$ for
        each known zero~$t_k$. For $\beta \neq 1/2$, $|L(\beta, t)|$ remains bounded
        away from zero for all $t \in [0, 1000]$.

  \item \textbf{System~2 ($G$):} $G(1/2, t)$ exhibits sign changes at exactly the same
        positions as the Hardy $Z$-function $Z_H(t)$, with the detected zero positions
        matching Odlyzko's tables~\cite{odlyzko1989} to within~$10^{-4}$. The zero count
        matches the Riemann--von Mangoldt prediction.

  \item \textbf{Rate function:} $F_1(1/2, t)$ exhibits sharp peaks (approaching~$1$, or
        $\ln 2$ under the Wei et~al.\ $\log$-base-$2$ normalization) at each zero, with
        $F_1(1/2, t) \approx 1/2$ (or $(1/2)\ln 2$ under Wei et~al.\ normalization)
        between zeros. The 2D phase diagram $|L(\beta, t)|$ in the $(\beta, t)$ plane
        shows dark lines (near-vanishing $|L|$) exclusively at $\beta = 1/2$.

  \item \textbf{Off-critical-line scan:} No values of $|L(\beta, t)|$ below the threshold
        $\eps = 10^{-6}$ for any $\beta \neq 1/2$ in the scanned range.
\end{enumerate}


\subsection{If RH Fails (Unexpected Scenario)}
\label{sec:rh-fails}

If a counterexample to~RH exists within our scanned range, we would observe:
\begin{enumerate}
  \item A value $(\beta_0, t_0)$ with $\beta_0 \neq 1/2$ where
        $|L(\beta_0, t_0)| < \eps$.
  \item Upon refinement, the zero would persist and sharpen ($|L|$ decreasing with
        increasing~$N$).
  \item The phase diagram would show an off-critical-line zero---a dark point or line at
        $\beta \neq 1/2$.
  \item Direct evaluation of $\zeta(\beta_0 + it_0)$ via independent methods (e.g.,
        Borwein acceleration or Euler--Maclaurin) would confirm the
        zero.\footnote{System~2's $G$ cannot be used for off-line confirmation due to the
        false-positive issue described in Section~\ref{sec:system2}.}
\end{enumerate}

Such a discovery would require extraordinary care in verification: independent
reimplementation, arbitrary-precision computation, and comparison with direct zeta
evaluation. The burden of proof for disproving~RH is immense, and any claimed
counterexample would need to withstand intense scrutiny.


\subsection{Limitations}
\label{sec:limitations}

We emphasize the inherent limitations of the numerical approach:
\begin{enumerate}
  \item \textbf{Cannot prove~RH.} Even scanning up to $T = 10^{10}$ and finding all zeros
        on the critical line does not prove that all zeros (of which there are infinitely
        many) lie on the critical line. Proof requires a fundamentally different
        (analytical or algebraic) approach.

  \item \textbf{Finite precision (FP64 wall).} Our computations use IEEE~754
        double-precision arithmetic ($\sim 15.9$ decimal digits). This imposes two
        distinct limitations: (a)~a zero at $\beta_0 = 1/2 + 10^{-20}$ would be
        indistinguishable from $\beta_0 = 1/2$; (b)~for $t > 10^7$, trigonometric
        argument reduction degrades precision below $8$ significant digits, rendering
        zero detection unreliable (see Section~\ref{sec:trig} for the full analysis).
        Extending the implementation to arbitrary-precision arithmetic
        (MPFR\footnote{The GNU MPFR library provides correctly-rounded multiple-precision
        floating-point computation with arbitrary mantissa width, at the cost of
        $O(p\log p)$ per operation for $p$-bit precision. For $p = 128$ (quad precision),
        the overhead is approximately $4$--$10\times$ relative to FP64; for $p = 256$,
        approximately $10$--$25\times$. This is a tractable overhead for extending
        reliable computation to $t \sim 10^{15}$ or beyond.}) would extend the reliable
        range to arbitrarily large~$t$, at the cost of proportionally increased
        computation time.

  \item \textbf{Finite scan range.} Zeros with $\im(s) > T_{\max}$ are not tested.
        Counterexamples to~RH, if they exist, might occur only at astronomically large
        heights. The first known failure of Gram's law occurs at Gram point $g_{126}$,
        and it is conceivable that the first off-critical-line zero (if any) could be at
        heights far beyond computational reach.

  \item \textbf{Indirect detection.} We detect zeros of~$L$ and~$G$, which are related
        to but not identical with~$\zeta$. Any discrepancy between $L$-zeros and
        $\zeta$-zeros must be attributed to finite-$N$ effects rather than interpreted as
        a new mathematical result.
\end{enumerate}

%% ====================================================================
\section{Relation to Prior Work}
\label{sec:prior}
%% ====================================================================

\subsection{Odlyzko's Numerical Verification}
\label{sec:odlyzko}

Andrew Odlyzko's seminal computations~\cite{odlyzko1989,odlyzko2001} focused on
computing specific high zeros of~$\zeta$---around the $10^{20}$th and $10^{22}$nd
zeros---and large blocks of consecutive zeros near these heights, providing crucial
statistical data on zero spacing and GUE pair correlation. The systematic verification
that the first $10^{13}$ nontrivial zeros of~$\zeta$ lie on the critical line was carried
out by Gourdon~\cite{gourdon2004}. Both Odlyzko's and Gourdon's methods use the
Odlyzko--Sch\"onhage algorithm, which evaluates $\zeta(1/2 + it)$ at many points
simultaneously using the Riemann--Siegel formula and FFT-based polynomial evaluation,
achieving $O(T^{1/2+\eps})$ complexity for all zeros up to height~$T$. These remain the
gold standard for numerical verification of~RH.

Our work does not compete with Odlyzko's in terms of the number of zeros verified.
Instead, it provides:
\begin{itemize}
  \item An alternative computational framework based on quantum physics.
  \item A cross-check using an entirely different mathematical representation (eta
        function and $Z$-function via DQPT order parameters rather than direct zeta
        evaluation).
  \item Visualization of the $\beta$-dependence (the 2D phase diagram), which is not part
        of the standard numerical verification workflow.
\end{itemize}


\subsection{Berry--Keating Conjecture}
\label{sec:berry-keating}

Berry and Keating~\cite{berrykeating1999} conjectured that the Hilbert--P\'olya operator
for~RH is related to the quantization of the classical Hamiltonian $H = xp$. Specifically,
they proposed that the operator $(xp + px)/2$, suitably regularized, has eigenvalues
related to the nontrivial zeta zeros. This connects~RH to quantum chaos: the zeta zeros
would be energy levels of a chaotic quantum system, and their statistics (Montgomery's
pair correlation conjecture~\cite{montgomery1973}, confirmed numerically by Odlyzko)
match those of random matrix theory (GUE ensemble).

The Wei et~al.\ construction~\cite{wei2025} is distinct from Berry--Keating in that it
does not seek an operator whose \emph{spectrum} encodes the zeros, but rather constructs
systems whose \emph{dynamical phase transitions} (temporal non-analyticities) correspond
to zeros. The spectrum of $H_0 = \sum\ln(n)\,|n\rangle\langle n|$ is logarithmic by
design, and the zeros emerge in the time domain rather than the energy domain.


\subsection{Connes' Operator-Algebraic Approach}
\label{sec:connes}

Alain Connes~\cite{connes1999} formulated a trace formula relating the zeros of~$\zeta$
to the spectrum of a certain operator on an adelic space, providing a noncommutative
geometry framework for~RH. This approach has led to deep structural insights but has not
yet yielded a proof of~RH. Our computational work is complementary: it provides numerical
data that any future proof (whether through Connes' framework or otherwise) must be
consistent with.


\subsection{Guth--Maynard 2024 Breakthrough}
\label{sec:guth-maynard}

In 2024, Larry Guth and James Maynard~\cite{guthmaynard2024} proved improved large-value
estimates for Dirichlet polynomials, which imply that at most $O(T^{3/5+\eps})$ zeros of
$\zeta(s)$ with $0 < \im(s) \le T$ can lie off the critical line, improving the previous
$O(T^{3/4+\eps})$ bound (Ingham~\cite{ingham1940}; refined by Huxley and others). Note
that the qualitative conclusion---that $100\%$ of zeros lie on the critical line in a
density sense (i.e., $N_{\mathrm{off}}(T)/N(T) \to 0$)---already followed from the prior
$O(T^{3/4+\eps})$ bound since $T^{3/4}/N(T) \to 0$. Guth--Maynard's contribution is the
improved quantitative exponent, which represents the first improvement to the zero density
estimate in the critical strip in over $80$~years.

Our numerical investigation probes the complementary question: among the specific zeros
in our scanned range, are there any off the critical line? Guth--Maynard tells us that
at most $T^{3/5+\eps}$ could be off-line up to height~$T$; our computation would detect
any such zeros within $[0, T_{\max}]$.


\subsection{How Our Implementation Differs}
\label{sec:novelty}

Our implementation is, to our knowledge, the first public classical simulation of the
Wei et~al.\ DQPT framework~\cite{wei2025}. It differs from all the above in:
\begin{enumerate}
  \item \textbf{Physical framing.} We compute quantum-mechanical observables (coherence
        functions, Loschmidt echoes, rate functions) rather than directly evaluating
        $\zeta$ or $L$-functions.
  \item \textbf{2D scanning.} We scan the full $(\beta, t)$ plane, not just the critical
        line.
  \item \textbf{DQPT characterization.} We identify zeros through their physical
        manifestation as phase transitions (singularities in the rate function), not
        merely as zero-crossings.
  \item \textbf{Dual systems.} We implement both System~1 (eta-based) and System~2
        ($Z$-function-based), providing internal cross-validation.
\end{enumerate}

%% ====================================================================
\section{Directions Toward Proof}
\label{sec:directions}
%% ====================================================================

As discussed in Section~\ref{sec:scope}, the DQPT--RH equivalence is a reformulation
rather than a proof mechanism. Nevertheless, the physical language suggests several
directions that might, in principle, lead to analytical progress on~RH. We outline these
directions here, emphasizing that none currently constitutes a proof strategy; they are
research programs that could benefit from the physical intuition provided by the DQPT
framework.


\subsection{Lee--Yang Mechanism for Alternating Series}
\label{sec:lee-yang-mech}

The Lee--Yang theorem~\cite{leeyang1952} states that for certain classes of partition
functions (e.g., ferromagnetic Ising models), all zeros of the partition function $Z(z)$
as a function of the fugacity~$z$ lie on the unit circle $|z| = 1$. The DQPT framework
interprets the eta function as a partition function of a system with logarithmic spectrum
and alternating (anti-ferromagnetic) couplings.

If one could establish a Lee--Yang-type theorem for the partition function
\begin{equation}\label{eq:ly-eta}
  Z(s) = \sum_{n=1}^{\infty} (-1)^{n+1}\,n^{-s},
\end{equation}
proving that its zeros (in the $s$-plane) are confined to the line $\re(s) = 1/2$, this
would prove~RH. The key challenge is that the Lee--Yang theorem in its classical form
applies to polynomial partition functions with specific positivity conditions, which do
not directly apply to the Dirichlet eta function. Extending Lee--Yang-type results to
infinite series with complex coefficients is a major open problem in mathematical
physics.


\subsection{Energy--Entropy Balance at $\beta = 1/2$}
\label{sec:energy-entropy}

In statistical mechanics, phase transitions occur when the free energy $F = E - TS$
(energy minus temperature times entropy) has a non-analyticity. At the critical
temperature~$T_c$ (equivalently, critical inverse temperature~$\beta_c$), the energy and
entropy contributions to the free energy are precisely balanced.

In the DQPT framework, the critical inverse temperature is $\beta_c = 1/2$. The
``energy'' of the system is related to the Dirichlet series
$\sum n^{-s}\ln n$ (the derivative of the partition function), and the ``entropy'' is
related to the combinatorial structure of the alternating series. The conjecture that
$\beta_c = 1/2$ is the \emph{unique} critical point is equivalent to~RH.

A proof strategy along these lines would need to show that the energy--entropy balance
condition is satisfied only at $\beta = 1/2$, for all~$t$. This requires understanding
the analytic structure of the free energy in the complex $\beta$-plane, which brings the
problem back to the original analytic number theory.

\paragraph{The Kramers--Wannier duality analogy.}
Perhaps the most promising physical mechanism for fixing the critical point at
$\beta = 1/2$ comes from duality. In the two-dimensional Ising model, the
Kramers--Wannier duality maps the partition function at inverse temperature~$\beta$ to
the partition function at a dual temperature~$\beta^*$, with the relation
$\sinh(2\beta)\sinh(2\beta^*) = 1$. The critical point is the self-dual point
$\beta_c = \beta_c^*$, which uniquely determines~$T_c$.

The functional equation of the Riemann zeta function plays a formally analogous role:
\begin{equation}\label{eq:xi-duality}
  \xi(s) = \xi(1 - s).
\end{equation}
This is a duality that maps $s$ to $1 - s$, fixing the critical line $\re(s) = 1/2$ as
the self-dual locus. In the DQPT framework, this becomes: the rate function
$F_1(\beta, t)$ satisfies a duality $F_1(\beta, t) \longleftrightarrow F_1(1 - \beta, t')$
(where the mapping $t \to t'$ involves the Riemann--Siegel theta function). The critical
point $\beta_c = 1/2$ is the unique self-dual point, just as in the Ising model.

The key question is whether this duality can be promoted to a \emph{proof} that all DQPTs
occur at the self-dual point $\beta = 1/2$. In the Ising model, the Kramers--Wannier
duality combined with the uniqueness of the phase transition (which follows from the
Peierls argument) proves that the critical temperature is at the self-dual point. The
analogue for the DQPT framework would require:
\begin{enumerate}
  \item[(a)] Proving a ``uniqueness of the phase transition'' theorem: that for each~$t$,
             there is at most one value of $\beta \in (0, 1)$ where $F_1$ has a
             singularity.
  \item[(b)] Combining this with the duality $\xi(s) = \xi(1 - s)$ to conclude that this
             unique critical~$\beta$ must be $1/2$.
\end{enumerate}
Neither (a) nor (b) is established, but this duality-based approach represents the most
physically motivated path from the DQPT framework to a proof of~RH. This is, in our
view, the central open question raised by the Wei et~al.\ construction.


\subsection{Universality of the DQPT Critical Point}
\label{sec:universality}

In equilibrium statistical mechanics, universality refers to the phenomenon that critical
exponents and scaling functions are independent of microscopic details. If the DQPT at
$\beta = 1/2$ exhibits universality in a suitable sense---meaning that perturbations of
the logarithmic spectrum $E_n = \ln(n)$ do not move the critical point away from
$\beta = 1/2$---this would suggest a structural reason for~RH.

Specifically, consider deformed spectra $E_n = \ln(n) + \eps\,f(n)$ for small~$\eps$ and
various perturbations~$f$. If the DQPT critical points remain at $\beta = 1/2$ for all
sufficiently small perturbations (structural stability), this would provide evidence that
$\beta = 1/2$ is ``locked in'' by some symmetry or topological constraint. Investigating
the stability of the DQPT under spectral perturbations is a concrete computational
project that could be carried out within our framework.


\subsection{Selberg Trace Formula Connection}
\label{sec:selberg}

The Selberg trace formula (Selberg, 1956) relates the spectrum of the Laplace--Beltrami
operator on a compact hyperbolic surface $\Gamma\backslash\mathbb{H}$ to the lengths of
closed geodesics on the surface. The Selberg zeta function $Z_\Gamma(s)$ is an entire
function whose zeros encode both the eigenvalues of the Laplacian and the topology of
the surface.

The analogy with the Riemann zeta function is deep: the Selberg zeta function satisfies a
functional equation, has a Riemann Hypothesis analog (which is \emph{proved} for compact
surfaces), and its zeros are related to a self-adjoint operator (the Laplacian). If the
Wei et~al.\ DQPT construction could be realized on a hyperbolic surface---with the
logarithmic spectrum $E_n = \ln(n)$ arising naturally from the geodesic lengths---then the
proven Selberg~RH might provide leverage toward the classical~RH.

The key obstacle is that the Selberg zeta function is a product over \emph{primitive}
closed geodesics, while the Riemann zeta function is a product over \emph{primes}.
Establishing a precise correspondence between geodesic lengths and prime numbers (beyond
the formal analogy) remains a central open problem.


\subsection{Random Matrix Universality from the DQPT Perspective}
\label{sec:rmt}

Montgomery's pair correlation conjecture~\cite{montgomery1973}, supported by Odlyzko's
numerical evidence~\cite{odlyzko1989}, asserts that the local statistics of zeta zeros
match those of eigenvalues of random matrices from the Gaussian Unitary Ensemble (GUE).
In the DQPT framework, this would correspond to the statement that the spacing statistics
of DQPT times~$t_k$ follow GUE statistics.

\paragraph{Integrability obstruction.}
A crucial limitation must be acknowledged: the Hamiltonian
$H = \sigma_z \otimes H_0$ with diagonal $H_0 = \sum \ln(n)\,|n\rangle\langle n|$
defines a prototypically \emph{integrable} system (a Gaudin-magnet-type central spin
model). The eigenstates are product states of the probe spin and bath energy eigenstates,
with no level repulsion. Integrable quantum systems generically exhibit \emph{Poisson}
level statistics (uncorrelated eigenvalue spacings), not GUE statistics. Therefore, the
DQPT quantum system described in this paper \emph{cannot} provide a physical mechanism
for the observed GUE statistics of zeta zeros through its own dynamics.

The GUE universality of the zeta zeros must instead arise from the number-theoretic
structure of the Riemann zeta function itself (encoded in the logarithmic spectrum
$E_n = \ln n$ and the resulting Dirichlet series), not from the quantum dynamics of the
probe-bath coupling. Any connection between the DQPT framework and random matrix theory
would need to operate at the level of the mathematical structure of the partition
function, not through a quantum chaos mechanism in the physical
system~\cite{berrykeating1999}. This remains an open problem.

%% ====================================================================
\section{Implementation Architecture}
\label{sec:architecture}
%% ====================================================================

\subsection{File Structure}
\label{sec:files}

The project is organized into four directories: \texttt{src/} (core computation),
\texttt{test/} (validation), \texttt{gpu/} (CUDA kernels), and \texttt{data/}
(reference tables).

\paragraph{Core source files (\texttt{src/}).}

\begin{longtable}{@{}p{4.2cm}p{8.5cm}@{}}
\toprule
\textbf{File} & \textbf{Purpose} \\
\midrule
\endhead
\texttt{dirichlet.h/.cpp}      & Partial Dirichlet series $S_N(s)$ and $\eta_N(s)$ \\
\texttt{partition.h/.cpp}      & Partition function $Z(\beta) = \sum n^{-\beta}$ with precomputation \\
\texttt{theta.h/.cpp}          & Riemann-Siegel theta function $\vartheta(t)$; Stirling expansion with error control \\
\texttt{system1.h/.cpp}        & $\mathcal{L}(\beta,t)$: accumulated phase factor (System~1) with alternating series \\
\texttt{system2.h/.cpp}        & $\mathcal{G}(\beta,t)$: Loschmidt amplitude (System~2) with $Z$-function connection \\
\texttt{rate\_function.h/.cpp} & Rate function $\mathcal{F}(\beta,t)$ with underflow guard and sentinel detection \\
\texttt{zeros.h/.cpp}          & Zero detection: bisection, Gram points, Turing zero-count check \\
\texttt{euler\_maclaurin.h/.cpp} & Euler-Maclaurin tail correction; Bernoulli number computation \\
\texttt{kahan.h}               & Kahan compensated summation \\
\texttt{main.cpp}              & Entry point; CLI argument parsing \\
\texttt{config.h}              & Compile-time and runtime configuration \\
\bottomrule
\end{longtable}

\paragraph{Test suite (\texttt{test/}).}

\begin{longtable}{@{}p{5cm}p{7.7cm}@{}}
\toprule
\textbf{File} & \textbf{Purpose} \\
\midrule
\endhead
\texttt{test\_dirichlet.cpp}        & Unit tests for $S_N$ and $\eta_N$ against \texttt{mpmath} \\
\texttt{test\_theta.cpp}            & Unit tests for $\vartheta(t)$ \\
\texttt{test\_known\_zeros.cpp}     & Verification against first 100 known zeros \\
\texttt{test\_convergence.cpp}      & Convergence rate $O(N^{-\beta})$ verification \\
\texttt{test\_off\_line.cpp}        & Off-critical-line scan \\
\texttt{test\_rate\_function.cpp}   & Rate function value checks (Controls 1--2) \\
\texttt{test\_cross\_validation.cpp}& System~1 vs System~2 consistency \\
\texttt{reference\_values.h}        & High-precision reference values from \texttt{mpmath} \\
\bottomrule
\end{longtable}

\paragraph{GPU kernels (\texttt{gpu/}).}

\begin{longtable}{@{}p{4.2cm}p{8.5cm}@{}}
\toprule
\textbf{File} & \textbf{Purpose} \\
\midrule
\endhead
\texttt{sweep.cu}           & CUDA kernel for $(\beta,t)$ grid evaluation \\
\texttt{gpu\_dirichlet.cuh} & Device-side Dirichlet series \\
\texttt{gpu\_config.cuh}    & GPU-specific configuration \\
\bottomrule
\end{longtable}

\paragraph{Data and output.}
\texttt{data/zeros\_known.txt} contains the first $100+$ known nontrivial zeta zeros
(sourced from the LMFDB~\cite{lmfdb} and Odlyzko~\cite{odlyzko1989}).
The \texttt{results/} directory is populated at runtime.


\subsection{Language Choices}
\label{sec:languages}

\paragraph{Core computation: C++17.}
Chosen for:
\begin{itemize}
  \item Direct control over floating-point arithmetic (IEEE~754 compliance via compiler
        flags \texttt{-ffp-contract=off}, \texttt{-fno-fast-math}).
  \item Efficient memory management for large arrays.
  \item Compatibility with CUDA (which requires C++ host code).
  \item No garbage collection pauses during long computations.
\end{itemize}

\paragraph{GPU acceleration: CUDA C++.}
Chosen for:
\begin{itemize}
  \item Native support on the target hardware (NVIDIA RTX~2060).
  \item Fine-grained control over thread/block organization.
  \item Access to hardware special function units for \texttt{sin}/\texttt{cos}/\texttt{exp}.
\end{itemize}

\paragraph{Test validation: Python/mpmath (external).}
Reference values are precomputed using \texttt{mpmath} with $50$-digit precision and
hardcoded into the C++ test suite. This removes any runtime dependency on Python.


\subsection{Dependencies}
\label{sec:deps}

The implementation is designed to minimize external dependencies:
\begin{itemize}
  \item \textbf{C++ standard library} (C++17): containers, algorithms, I/O,
        \texttt{<cmath>}.
  \item \textbf{CUDA Toolkit} ($\ge 11.0$): for GPU kernels. Optional; the code compiles
        and runs in CPU-only mode without CUDA.
  \item \textbf{No other external dependencies.} No Boost, no Eigen, no BLAS/LAPACK, no
        numerical libraries. All required mathematical functions (Stirling expansion,
        Kahan summation, Bernoulli numbers, bisection) are implemented from scratch for
        full transparency and auditability.
\end{itemize}


\subsection{Build System}
\label{sec:build}

A single Makefile with the following targets:

\begin{center}
\begin{tabular}{@{}ll@{}}
\toprule
\textbf{Command} & \textbf{Action} \\
\midrule
\texttt{make}       & Build CPU-only version \\
\texttt{make gpu}   & Build with CUDA GPU support \\
\texttt{make test}  & Build and run all unit tests \\
\texttt{make clean} & Remove build artifacts \\
\bottomrule
\end{tabular}
\end{center}

Compiler flags:
\begin{itemize}
  \item \texttt{-std=c++17 -O2 -Wall -Wextra -Wpedantic}
  \item \texttt{-ffp-contract=off} (prevent fused multiply-add, which can change results)
  \item \texttt{-fno-fast-math} (ensure IEEE~754 compliance)
  \item For CUDA: \texttt{-arch=sm\_75} (Turing architecture, RTX~2060)
\end{itemize}

%% ====================================================================
\section{Success Criteria (Quantitative)}
\label{sec:criteria}
%% ====================================================================

Each criterion is stated as a specific numerical test with a pass/fail threshold.


\subsection{Verification Criteria (Must Pass)}
\label{sec:verification}

\paragraph{V1. Known zero recovery (System~1).}
For each of the first $100$ known nontrivial zeros $t_k$ ($k = 1, \ldots, 100$):
\begin{equation}\label{eq:V1}
  |L(0.5, t_k)| < 10^{-4} \qquad \text{with } N = 10^4
  \text{ and Euler-Maclaurin correction.}
\end{equation}
Pass: all~$100$. Fail: any one exceeds threshold.

\paragraph{V2. Known zero recovery (System~2).}
For each of the first $100$ known zeros:
\begin{equation}\label{eq:V2}
  G(0.5, t_k) \text{ changes sign in an interval of width } < 10^{-3}
  \text{ centered at } t_k.
\end{equation}
Pass: all $100$ sign changes detected within tolerance. Fail: any missed or mislocated.

\paragraph{V3. Zero count accuracy.}
For $T \in \{100, 500, 1000\}$:
\begin{equation}\label{eq:V3}
  |N_{\mathrm{detected}}(T) - N_{\mathrm{predicted}}(T)| \le 1,
\end{equation}
where $N_{\mathrm{predicted}}(T) = \mathrm{round}((1/\pi)\,\vartheta(T) + 1)$. The
tolerance of~$1$ accounts for the $S(T) = (1/\pi)\arg\zeta(1/2 + iT)$ fluctuation term,
which can cause the actual zero count to differ from the smooth prediction by $\pm 1$ at
specific $T$ values.

Pass: agreement within tolerance for all three $T$ values. Fail: discrepancy $> 1$ for
any~$T$.

\paragraph{V4. Rate function at zeros.}
For the first $10$ known zeros:
\begin{equation}\label{eq:V4}
  |F_1(0.5, t_k) - 1.0| < 0.1
\end{equation}
(natural log normalization; equivalent to $|F_1 - \ln 2| < 0.07$ under Wei et~al.\
$\log$-base-$2$ normalization).

Pass: all $10$ within tolerance. Fail: any one outside.

\paragraph{V5. Rate function away from zeros.}
For $t = 50.0$ (not near any zero) and $\beta \in \{0.3, 0.5, 0.7\}$:
\begin{equation}\label{eq:V5}
  |F_1(\beta, 50.0) - (1 - \beta)| < 0.05
\end{equation}
(natural log normalization; equivalent to $|F_1 - (1 - \beta)\ln 2|$ under Wei et~al.\
normalization).

Pass: all three $\beta$ values within tolerance. Fail: any one outside.

\paragraph{V6. Cross-validation (System~1 vs System~2).}
For all detected zeros, the positions from System~1 and System~2 agree:
\begin{equation}\label{eq:V6}
  |t_k^{(1)} - t_k^{(2)}| < 10^{-3}.
\end{equation}
Pass: all zeros agree. Fail: any discrepancy exceeding threshold.

\paragraph{V7. Convergence rate.}
Two convergence tests are performed:
\begin{enumerate}
  \item[(a)] \emph{At a known zero:} $s = 0.5 + 14.134725i$ (first zeta zero). The
             sequence $|\eta_N(s)|$ for $N = 100, 1000, 10000$ must decrease as
             $O(N^{-0.5})$. The ratio $|\eta_{10N}(s)|/|\eta_N(s)|$ should be
             approximately $1/\sqrt{10} \approx 0.316$.

             Pass: ratio in $[0.25, 0.40]$. Fail: ratio outside $[0.20, 0.50]$.

  \item[(b)] \emph{At a non-zero point:} $s = 0.5 + 50.0i$ (not a zeta zero;
             $\eta(s) \approx -0.85 + 0.36i$). The sequence $\eta_N(s)$ must converge to
             the reference value $\eta(s)$ (computed via Borwein acceleration). The
             relative error $|\eta_N(s) - \eta(s)|/|\eta(s)|$ for $N = 100, 1000, 10000$
             must decrease as $O(N^{-0.5})$.

             Pass: relative error at $N = 10000$ is below $10^{-2}$. Fail: relative
             error above $10^{-1}$.
\end{enumerate}


\subsection{RH-Relevant Criteria}
\label{sec:rh-criteria}

\paragraph{R1. Off-critical-line absence of zeros.}
For all $\beta \in [0.1, 0.9]$ with spacing $0.005$, excluding $[0.4975, 0.5025]$, and
$t \in [0, 1000]$:
\begin{equation}\label{eq:R1}
  \min_{t}\,|L(\beta, t)| > 10^{-3}.
\end{equation}
Pass: no $\beta$ shows a near-zero. Fail: any $\beta$ shows $|L| < 10^{-3}$ (triggers
Stage~4 investigation, but is not by itself a failure of the implementation).

\paragraph{R2. Phase diagram confinement.}
The 2D heat map of $-\log|L(\beta, t)|$ in the $(\beta, t)$ plane shows strong signal
(large $-\log|L|$) exclusively in the column $\beta = 0.5$, with no comparable signal at
other $\beta$ values.

Pass: visual and quantitative confirmation. Fail: comparable signal at $\beta \neq 0.5$.

%% ====================================================================
\section{Timeline and Milestones}
\label{sec:timeline}
%% ====================================================================

\subsection*{Phase~1: Core Functions and Unit Tests (Weeks 1--2)}

\paragraph{Deliverables.}
\begin{itemize}
  \item Implemented and tested: $S_N(s)$, $\eta_N(s)$, $Z(\beta)$, $\vartheta(t)$,
        Kahan summation, Euler--Maclaurin correction.
  \item Unit test suite passing against \texttt{mpmath} reference values.
  \item Verified: $10$-digit agreement for $S_N$, $\vartheta$, $\eta$ at selected test
        points.
\end{itemize}

\paragraph{Milestone.}
\texttt{make test} passes all Phase~1 tests.


\subsection*{Phase~2: System Implementation and Known-Zero Verification (Weeks 3--4)}

\paragraph{Deliverables.}
\begin{itemize}
  \item Implemented: $L(\beta, t)$ (System~1), $G(\beta, t)$ (System~2), $F(\beta, t)$
        (rate function).
  \item Zero detection algorithm (bisection $+$ Gram points $+$ Turing check).
  \item Verification: $L$ and $G$ vanish at first $100$ known zeros (criteria V1, V2).
  \item Zero count matches Riemann--von Mangoldt (criterion V3).
  \item Rate function values correct (criteria V4, V5).
  \item Cross-validation passes (criterion V6).
  \item Convergence rate verified (criterion V7).
\end{itemize}

\paragraph{Milestone.}
All verification criteria (V1--V7) pass.


\subsection*{Phase~3: Full $\beta$-$t$ Scan (Weeks 5--6)}

\paragraph{Deliverables.}
\begin{itemize}
  \item CPU-based scan of the $(\beta, t)$ plane with $\Delta\beta = 0.005$,
        $\Delta t = 0.1$, $T_{\max} = 1000$.
  \item Off-critical-line analysis (criterion R1).
  \item 2D phase diagram generation and analysis (criterion R2).
  \item Investigation and resolution of any flagged candidates from Stage~3/4 of the
        test protocol.
\end{itemize}

\paragraph{Milestone.}
Criteria R1 and R2 evaluated. No false off-line zeros after refinement.


\subsection*{Phase~4: GPU Acceleration and Large-Scale Computation (Weeks 7--10)}

\paragraph{Deliverables.}
\begin{itemize}
  \item CUDA kernel for $(\beta, t)$ grid evaluation.
  \item Performance benchmarking on RTX~2060.
  \item Extended scan: $T_{\max} = 10{,}000$ (capturing $\sim 10{,}142$ zeros).
  \item Fine $\beta$ grid: $\Delta\beta = 0.001$ for refined scans near flagged
        candidates.
  \item High-resolution phase diagram.
\end{itemize}

\paragraph{Milestone.}
GPU scan completes in $< 1$~hour for $T_{\max} = 10{,}000$. All zeros confirmed on
critical line.


\subsection*{Phase~5: Results Analysis and Documentation (Weeks 11--12)}

\paragraph{Deliverables.}
\begin{itemize}
  \item Summary of all detected zeros and their $(\beta, t)$ coordinates.
  \item Statistical analysis: distribution of $|L(\beta, t)|$ as a function of~$\beta$.
  \item Phase diagrams at multiple resolutions.
  \item Comparison with published tables (number of matching zeros, position accuracy).
  \item Final report documenting results, methodology validation, and conclusions.
\end{itemize}

\paragraph{Milestone.}
All success criteria met. Results documented.


\addcontentsline{toc}{section}{References}
\bibliographystyle{unsrtnat}
\bibliography{references}

\end{document}
